%% 第26回編集委員会（物理第4巻）：組 電W4
%% 第27回 確認の直し（修4-3。E3-1〜E3-30・G-1〜G-15）：記録は tools/FIX-phys4-review-3.md、検算は checks/物理第4巻/確認/修4-3/check_fix4_3.py
%% p4-07.tex  電-7　鉄と磁石：磁性体と磁気回路（★1〜5・3段）
%% 中身は計画書 tools/phys4-plan.md の B「電-7」。公式の持ち主は A（A2 の磁化・A41 の物質の B＝μH・A72）。
%%   H を「磁場の強さ」としてこの章で出す（Y-2）。B は電-6 と同じく「磁場（磁束密度）」。
%%   磁区・交換相互作用・キュリー温度のしくみは「物理シリーズの『原子』（原-7・原-8）で学ぶ」と書いて結果だけ使った（Y-10）。
%%   磁気回路（⑤）は、前付の「相似の地図」の表に行を足す題材（計画書 0 の5）。⑤の表 ph:電-7:tab-souji に、電気回路と磁気回路の
%%   対応を作った。磁気抵抗は熱を出さずエネルギーをためる（電気抵抗とは役わりが違う）ことを「せいかくに言うと」に書いた。
%%   磁束 Φ＝BS は電-8 ③が持ち主（A41）。この章では磁気回路に要るので⑤で量として出し、「電-8 ③で詳しく」と書いた
%%   （記号カードは置かない。電-8 と照らし合わせる：報告に書いた）。電磁石の吸引力 B²S/(2μ0) は電-8 ⑩が持ち主で、例題11 は結果を使った。
%%   ⑧のヒステリシス損の「BH 曲線の面積」は、電源がする仕事を電-8 ③（ファラデーの法則）の結果で書いた（式でわかる段）。
%%   電-1〜5・8 はほかの組が同時に書いているので、計画書 B の節の番号で指した（電-3 ⑨電場の境界条件、電-4 ③R＝ρl/S、
%%   電-4 ⑦キルヒホッフ、電-8 ③磁束とファラデー・⑩吸引力・⑪渦電流と積層鉄心）。
%% 数学との重なり（Y-9）：5-12（積分は面積）を⑧で、8-11（小さな袋と道）を⑦で使った。
%% 数値はすべて Python で計算した：checks/物理第4巻/電-7/calc_p4_07.py（出力 calc_p4_07_out.txt）。⑧の表計算の表も同じスクリプト。
%%   常用漢字の確かめ：checks/物理第4巻/電-6/kanji_p4_06_07.py
%% 記号（この章）：磁場（磁束密度）B、磁場の強さ H、磁化 M（電-8 の相互インダクタンス M とは別の章）、透磁率 μ・比透磁率 μ_r、
%%   磁化率 χ、巻数 N・1 m あたりの巻数 n、電流 I、起磁力 NI（記号 𝓕 は専門書の囲みだけ）、磁束 Φ、磁気抵抗 𝓡_m
%%   （電気抵抗 R と書体と添字で分ける）、鉄心の長さ l・断面積 S、すき間の長さ l_g、磁石の長さ l_m、保磁力 H_c、残留磁束密度 B_r、
%%   リコイル比透磁率 μ_rec、起電力 𝓔（相似の表の電気の列：Y-3 の場の章の書き方）、密度 ρ、重力加速度 g（例題10）。
%%   磁気モーメントは記号を作らず「NIS」（電-6 ⑪）と書いた。
%% 出典（書くときに確かめた。11月5日）
%%   ・μ0：電-6 の頭のコメントと同じ（CODATA 2022、BIPM）。
%%   ・水の磁化率 −9.0×10⁻⁶（SI、体積あたり）と、水を浮かせる B dB/dz≈1.4×10³ T²/m：arXiv:2004.09683（超伝導線を使った反磁性浮上の論文。
%%       「water (χ = −9.0×10⁻⁶, ρ = 1000 kg m⁻³), B dB/dz = −1360 T² m⁻¹」）、Nottingham 大学の New J. Phys. の論文（反磁性浮上、1360 T²/m）、
%%       Leicester 大学の学生論文誌（1400 T²/m）。16 T でカエルを浮かせた実験：M. V. Berry and A. K. Geim, "Of flying frogs and levitrons",
%%       Eur. J. Phys. 18, 307–313 (1997)（書誌と 16 T は WebSearch の検索の抜粋で、複数の資料で照らし合わせた。原文は作業環境から届かず）。
%%   ・アルミニウムの磁化率 +2.2×10⁻⁵、銅 約 −1.0×10⁻⁵（SI、体積あたり）：Encyclopedia Magnetica の aluminium の項、
%%       テキサス大学 Fitzpatrick の電磁気の講義ノート（CRC Handbook を引く）の検索の抜粋（WebSearch。どちらも作業環境からは本文を読めなかった）。
%%   ・鉄のキュリー温度 770 ℃（1043 K）：ミシガン大学 材料科学の演示実験のページ（mse.engin.umich.edu "Curie Point of Iron"）、
%%       NETZSCH の用語集（WebSearch で照らし合わせた）。
%%   ・ネオジム磁石の残留磁束密度 約1.2〜1.4 T、フェライト磁石 約0.4 T：MISUMI の技術資料（jp.misumi-ec.com/pdf/fa/2018/p2_321.pdf：
%%       強力ネオジム 1.42 T 以上、ネオジム 1.26 T 以上、フェライト 0.385 T 以上）、ネオマグのカタログ（neomag.jp）。例題では 1.3 T・0.40 T と置いた。
%%       リコイル比透磁率（ネオジム 1.05、フェライト 1.1）は例題の仮の値（本文に「〜とします」と書いた）。
%%   ・電磁鋼板の飽和の目安 約2 T：日本製鉄の電磁鋼板の製品の説明（nipponsteel.com）、方向性電磁鋼板の特許（JPH0688174A：製品の磁束密度 1.9〜2 T）。
%%   ・無方向性電磁鋼帯 35A300 の鉄損 W15/50 が 3.00 W/kg 以下（記号の数が鉄損の上限）：JIS C 2552:2014 の規格票のプレビュー
%%       （webdesk.jsa.or.jp）と、規格の対応表の資料（WebSearch で照らし合わせた）。電磁鋼板の密度 7650 kg/m³ は例題の仮の値。
%%   ・最後の地磁気の逆転（約77万年前）と千葉セクション・チバニアン（2020年1月17日、国際地質科学連合の理事会で決定）：
%%       産業技術総合研究所 地質調査総合センター GSJ 地質ニュース（gsj.jp/data/gcn/gsj_cn_vol9.no8_p.231-232.pdf）、日本ジオパークネットワークの
%%       お知らせ（geopark.jp/jgn/news/2020/pdf/20200117.pdf）、千葉日報（2020年1月）で照らし合わせた（WebSearch）。
%%   ・地磁気（東京）の全磁力 約46000 nT：電-6 の頭のコメントと同じ（〔要確認〕：気象庁地磁気観測所・国土地理院に届かず）。
%% 〔要確認〕：鉄・電磁鋼板の比透磁率は材料と磁場の強さで大きく変わるので、本文は「数百〜数千」とし、例題の値は仮の値と書いた。
\makeatletter
% 棒磁石  \sxhmagnet{(中心)}{長さ}{太さ}{右の極の字}{左の極の字}
%   N極の側を赤、S極の側を青でぬる（極の字を入れかえても、赤＝N にそろう：E3-11）
\def\sxh@N{N}
\newcommand\sxhmagnet[5]{\begin{scope}[shift={#1}]
  \def\sxh@r{#4}\ifx\sxh@r\sxh@N\def\sxh@rc{bkRed!55}\def\sxh@lc{bkBlue!40}\else\def\sxh@rc{bkBlue!40}\def\sxh@lc{bkRed!55}\fi
  \filldraw[fill=\sxh@rc, draw=bkInk, line width=0.6pt] (0,{-#3/2}) rectangle ({#2/2},{#3/2});
  \filldraw[fill=\sxh@lc, draw=bkInk, line width=0.6pt] ({-#2/2},{-#3/2}) rectangle (0,{#3/2});
  \node[font=\sffamily\bfseries\footnotesize] at ({#2/4},0) {#4};
  \node[font=\sffamily\bfseries\footnotesize] at ({-#2/4},0) {#5};\end{scope}}
% 小さな磁石の矢印（磁区の中の向き）  \sxharr{(中心)}{角}{長さ}
\newcommand\sxharr[3]{\draw[bkRed!75!black, line width=0.9pt, -{Stealth[length=1.6mm]}] #1 ++({#2+180}:{#3/2}) -- ++(#2:#3);}
% 紙から手前へ・おくへ（電-6 と同じ形）
\newcommand\sxhout[2]{\begin{scope}[shift={#1}]\draw[bkInk, line width=0.6pt, fill=white] (0,0) circle (#2); \fill[bkInk] (0,0) circle ({0.32*#2});\end{scope}}
\newcommand\sxhin[2]{\begin{scope}[shift={#1}]\draw[bkInk, line width=0.6pt, fill=white] (0,0) circle (#2); \draw[bkInk, line width=0.6pt] (45:#2) -- (225:#2) (135:#2) -- (315:#2);\end{scope}}
\tikzset{sxh lab/.style={font=\sffamily\footnotesize, text=bkInk},
  sxh B/.style={bkBlue!80!black, line width=0.9pt, -{Stealth[length=2mm]}},
  sxh I/.style={bkOrange!85!black, line width=1.4pt, -{Stealth[length=2.4mm]}},
  sxh core/.style={fill=bkGray!35, draw=bkInk, line width=0.6pt, even odd rule}}
\makeatother

\phchapter{7}{鉄と磁石：\\磁性体と磁気回路}{磁石にくっついたくぎは、なぜ自分も磁石になる？}{\bklevelrange{1}{5}}{鉄と磁石：磁性体と磁気回路}
\bkmichishirube{5-12,8-11,電-3,電-4,電-6}{電-8,電-9,電-12,原-7,原-8}

\begin{goal}
\begin{itemize}[label={\color{bkTheme}●}, leftmargin=1.6\zw, itemsep=0pt, topsep=1pt]
\item \bklevel{1}\enspace 磁石に近づけた鉄が磁石になるわけを、小さな磁石の向きがそろう図で説明できる。
\item \bklevel{3}\enspace 磁化 $M$・磁場の強さ $H$・比透磁率を使って、鉄心の中の磁場を求められる。
\item \bklevel{3}\enspace 磁気回路を電気回路と同じ形で解き、すき間のある電磁石の磁束を求められる。
\item \bklevel{4}\enspace 鉄と空気の境目で磁力線がどう曲がるかを、境界条件から求められる。
\item \bklevel{5}\enspace $BH$ 曲線の面積からヒステリシス損を、減磁曲線から永久磁石の動作点を求められる。
\end{itemize}
\end{goal}

\phkoushiki{%
\phk{磁束密度と磁場の強さ $B=\mu H$（透磁率 $\mu$）}{高校物理}{ph:電-7:bmuh}
}

\begin{hatena}
\noindent\begin{minipage}[c]{0.55\linewidth}\raggedright
鉄のくぎの頭を、磁石にくっつけます。そのくぎの先に、別のくぎやクリップを近づけると、くっつきます。くぎが、磁石のように鉄を引きつけるようになったのです。

くぎは、もともと磁石ではありませんでした。磁石にくっついただけで、なぜ自分も磁石になるのでしょう。磁石から離すと、どうなるのでしょう。アルミの1円玉や10円玉は、磁石につきません。鉄と、ほかの金属は、何が違うのでしょう。
\end{minipage}\hfill
\begin{minipage}[c]{0.42\linewidth}\centering
\begin{tikzpicture}[x=1mm, y=1mm]
  \sxhmagnet{(24,42)}{30}{8}{N}{S}
  \fill[bkGray!70] (33.5,38) -- (34.5,38) -- (34.5,18) -- (34,15) -- (33.5,18) -- cycle;
  \fill[bkGray!80] (31.8,38) rectangle (36.2,39.4);
  \draw[bkGray!80!black, line width=1.2pt, rounded corners=1mm] (31,13.5) -- (31,5) -- (37,5) -- (37,11) -- (33,11) -- (33,7);
  \node[sxh lab, anchor=west] at (39,28) {くぎ};
  \node[sxh lab, anchor=west] at (39,8) {クリップ};
  \node[sxh lab, anchor=east, text=bkRed!80!black] at (29,10) {つく！};
\end{tikzpicture}
\end{minipage}
\end{hatena}

\noindent 電-6では、電流が磁場をつくることを学びました。では、電流の流れていない永久磁石は、なぜ磁場をつくるのでしょう。そして、鉄は、なぜ磁石に引きつけられ、自分も磁石になるのでしょう。この章では、物の中に入った磁場を調べます。鉄は、ソレノイドの中の磁場を何千倍にも強くします。モーター・変圧器・スピーカーには、必ず鉄や磁石が使われています。その設計に使う「磁気回路」は、電気回路とそっくりの形で解けます。

\begin{column}{磁石の石と、地球という磁石}
鉄を引きつける石（磁鉄鉱をふくむ天然の磁石）は、大昔から知られていました。この石を細長くけずって水に浮かべたり、糸でつるしたりすると、いつも同じ向きを向くことがわかり、方位磁針として船の旅に使われるようになりました。1600年、イギリスの医者ギルバートは『磁石について』という本を書きました。ギルバートは、磁石の石を球の形にけずった「小さな地球」の表面に小さな方位磁針を置き、磁針が場所によって決まった向きを向き、下へ傾くようすを調べました。そして、方位磁針が北を向くのは、地球そのものが大きな磁石だからだ、と考えました（電-6 ⑫の伏角）。磁石の性質を、実験で1つずつ確かめていった最初の本の1つです。
\end{column}

\begin{yattemiyou}[くぎを磁石にしよう]
\textbf{用意するもの}：鉄のくぎ（長さ\qty{5}{cm}くらい）、磁石（冷蔵庫の磁石より強い、棒や円板の磁石）、ゼムクリップ（10個くらい）、方位磁針\par\vspace{2pt}
\noindent\begin{minipage}[c]{0.64\linewidth}
\begin{enumerate}[label={\bkcircled[bkGreen]{\arabic*}}, leftmargin=2.4\zw, itemsep=1pt, topsep=0pt]
\item くぎの先を、クリップに近づける。つくか、つかないか。
\item 磁石のN極で、くぎの頭から先へ、同じ向きに10回こする（往復しない）。
\item くぎの先にクリップを何個つり下げられるか数える。くぎの先を方位磁針に近づけ、針のどちらの極が引きつけられるか調べる。
\item さらに10回こすって（合わせて20回）、③をする。
\item 次の日に、もう一度③をする。
\end{enumerate}
\end{minipage}\hfill
\begin{minipage}[c]{0.33\linewidth}\centering
\begin{tikzpicture}[x=1mm, y=1mm]
  \fill[bkGray!70] (2,10) -- (36,10) -- (40,11) -- (36,12) -- (2,12) -- cycle;
  \fill[bkGray!80] (0.6,9) rectangle (2,13);
  \sxhmagnet{(14,20)}{16}{6}{N}{S}
  \draw[bkGreen!50!black, line width=1pt, -{Stealth[length=2mm]}] (6,26) -- (32,26);
  \node[sxh lab, anchor=south] at (19,26.6) {同じ向きにこする};
  \node[sxh lab, anchor=north] at (20,8) {くぎ（頭 → 先）};
\end{tikzpicture}
\end{minipage}
\end{yattemiyou}

\begin{bkcenter}
\begin{bktable}{colspec={Q[c,m,wd=32mm] Q[c,m,wd=34mm] Q[c,m,wd=44mm]}, rowsep=1pt, row{1}={bg=bkTheme!10}}
こすった回数 & つり下げたクリップ & くぎの先が引きつけた針の極 \\
0回 & 0個 & どちらも少しだけ \\
10回 & 3個 & N極 \\
20回 & 5個 & N極 \\
20回（次の日） & 4個 & N極 \\
\end{bktable}
\bktabcaption{やってみようで、ミナさんが調べた結果}\label{ph:電-7:tab-yattemiyou}
\end{bkcenter}

\noindent やってみようの答え：表\ref{ph:電-7:tab-yattemiyou}のように、こする前のくぎは、クリップをつり下げられません（方位磁針の針は、どちらの極も少し引きつけます。針そのものが磁石なので、ただの鉄を引きつけるからです）。N極でこすると、くぎは磁石になり、くぎの先は針のN極を引きつけるようになります。くぎの先は\textbf{S極}になったのです（磁石のN極が最後に離れた所が S極になる）。こする回数をふやすと、少し強くなります。次の日には、少し弱くなっていますが、磁石のままです。くぎの中で何が起きたのかを、①〜③で見ていきます。

\begin{omoidasou}
\bkrecall{電-6}{電流がつくる磁場。長いソレノイドの中は $B=\mu_0nI$（$n$ は 1 m あたりの巻数）。高校の「磁場の強さ」$H=nI$ と $B=\mu_0H$（電-6 ③⑥）。}
\bkrecall{電-6}{円い電流は、平たい磁石のようにふるまう。コイルの「磁石としての強さ」は $NIS$（巻数×電流×面積）（電-6 ④⑪）。}
\bkrecall{電-4}{導線の抵抗は $R=\rho l/S$。長いほど大きく、太いほど小さい（電-4 ③）。キルヒホッフの法則（電-4 ⑦）。}
\bkrecall{5-12}{力×動いた距離のグラフの面積が仕事。グラフの面積は、細かく分けて台形の和で求められる（5-12②）。}
\bkrecall{8-11}{閉じた面から出る量を数える「袋」と、閉じた道にそってたす「道」（8-11⑦⑬）。}
\end{omoidasou}

\noindent この章の記号：磁化を $M$ と書きます（電-8 の相互インダクタンス $M$ とは別の量で、使う章が違います）。磁気抵抗は、電気抵抗 $R$ と区別するため、筆記体の $\mathcal R_{\mathrm m}$ と書きます。

%% ======================== 物理の道具 ========================
\bksection{doguu}{物理の道具}

\phstage{me}{磁石と鉄のようすを見る}{1}

\bksubsection{1}{磁石の性質：N極とS極}
\label{ph:電-7:sec1}
磁石には、鉄を強く引きつける所が2か所あり、これを\textbf{磁極}といいます。糸でつるした磁石の、北を向く磁極がN極、南を向く磁極がS極です。N極とS極は引き合い、N極どうし・S極どうしは退け合います（小学校の理科）。\bkidx{じきょく}{磁極}

磁石の上に紙をのせて砂鉄をまくと、砂鉄は磁力線の形（電-6 の図1）に並びます。砂鉄の1粒1粒が、磁石の磁場で小さな磁石になり、となりの粒とN極・S極でつながって、すじを作るのです。磁石を細かく切っていっても、どんなに小さなかけらにも N極と S極が必ずあります（図\ref{ph:電-7:fig-poles}の下）。磁石の性質は、とても小さな部分まで、ずっと続いているのです。

\begin{bkfigure}
\begin{tikzpicture}[x=1mm, y=1mm]
  % 引き合う
  \sxhmagnet{(0,0)}{20}{6}{N}{S}
  \sxhmagnet{(30,0)}{20}{6}{N}{S}
  \foreach \yy in {-2,0,2} {\draw[bkBlue!70!black, line width=0.6pt] (10,\yy) -- (20,\yy);}
  \draw[bkBlue!80!black, line width=0.8pt, -{Stealth[length=1.8mm]}] (14,0) -- (16.5,0);
  \node[sxh lab, anchor=north] at (15,-6) {N と S：引き合う};
  % 退け合う
  \begin{scope}[shift={(74,0)}]
  \sxhmagnet{(-15,0)}{20}{6}{N}{S}
  \sxhmagnet{(15,0)}{20}{6}{S}{N}
  \draw[bkBlue!70!black, line width=0.6pt] (-5,1.5) .. controls (-1,4) and (-1,10) .. (-3,14);
  \draw[bkBlue!70!black, line width=0.6pt] (5,1.5) .. controls (1,4) and (1,10) .. (3,14);
  \draw[bkBlue!70!black, line width=0.6pt] (-5,-1.5) .. controls (-1,-4) and (-1,-10) .. (-3,-14);
  \draw[bkBlue!70!black, line width=0.6pt] (5,-1.5) .. controls (1,-4) and (1,-10) .. (3,-14);
  \node[sxh lab, anchor=north] at (0,-15) {N と N：退け合う};
  \end{scope}
\end{tikzpicture}
\par\vspace{3mm}
\begin{tikzpicture}[x=1mm, y=1mm]
  \sxhmagnet{(0,0)}{40}{6}{N}{S}
  \draw[bkInk, line width=0.8pt, -{Stealth[length=2.4mm]}] (23,0) -- (31,0);
  \node[sxh lab, anchor=south] at (27,4) {2つに切る};
  \sxhmagnet{(44,0)}{20}{6}{N}{S}
  \sxhmagnet{(67,0)}{20}{6}{N}{S}
  \node[sxh lab, anchor=west, text width=30mm, align=left] at (79,0) {切った所に、新しい S極と N極ができる};
\end{tikzpicture}
\bkfigcaption{磁石の性質。上：磁極どうしの力と磁力線。下：磁石を切っても、N極だけ・S極だけにはならない。}\label{ph:電-7:fig-poles}
\end{bkfigure}

\bksubsection{2}{くぎが磁石になる：小さな磁石がそろう}
\label{ph:電-7:sec2}
どんどん細かくしていくと、最後には、鉄の原子1個1個が、小さな磁石になっています。鉄の原子の中の電子は、こまのように回る性質（\textbf{スピン}）をもち、そのために1つ1つが小さな磁石になっています（電-6 ④の円い電流が磁石になるのと似ていますが、スピンは量子の性質で、本当に回っているわけではありません）。\bkidx{すぴん}{スピン}

鉄では、となり合う原子の小さな磁石が、たがいに同じ向きにそろおうとする強い働きがあります。そのため、鉄の中は、小さな磁石の向きがそろった小部屋に分かれています。この小部屋を\textbf{磁区}といいます。\bkidx{じく}{磁区}1つの磁区の大きさは、1 mm の100分の1から10分の1くらい（材料で大きく違う）で、くぎ1本の中には、たくさんの磁区があります。

\begin{bkfigure}
\begin{tikzpicture}[x=1mm, y=1mm]
  % 左：ばらばら
  \draw[bkInk, line width=0.8pt, fill=bkGray!12] (0,0) rectangle (40,26);
  \draw[bkInk!70, line width=0.5pt] (14,0) -- (12,12) -- (0,14) (12,12) -- (22,15) -- (26,26) (22,15) -- (30,0) (22,15) -- (40,12);
  \sxharr{(6,6)}{0}{7}\sxharr{(6,20)}{90}{7}\sxharr{(21,6)}{180}{7}\sxharr{(15,21)}{-90}{7}\sxharr{(33,20)}{180}{7}\sxharr{(35,5)}{90}{7}
  \node[sxh lab, anchor=north] at (20,-1.5) {① 磁石に近づける前};
  \draw[bkInk, line width=0.9pt, -{Stealth[length=2.4mm]}] (44,13) -- (52,13);
  % 右：そろう
  \begin{scope}[shift={(56,0)}]
  \draw[bkInk, line width=0.8pt, fill=bkGray!12] (0,0) rectangle (40,26);
  \draw[bkInk!70, line width=0.5pt] (4,0) -- (3,8) -- (0,9) (3,8) -- (6,26) (6,26) (3,8) -- (40,6);
  \sxharr{(22,17)}{180}{9}\sxharr{(12,13)}{180}{9}\sxharr{(30,12)}{180}{9}\sxharr{(28,22)}{180}{9}\sxharr{(20,3)}{0}{6}\sxharr{(2,18)}{-90}{5}
  \node[sxh lab, anchor=north] at (20,-1.5) {② 磁石を近づけると};
  \sxhmagnet{(58,13)}{24}{7}{S}{N}
  \node[sxh lab, anchor=north, text width=30mm, align=center] at (58,8) {磁石のN極を\\くぎの右はしに};
  \end{scope}
\end{tikzpicture}
\bkfigcaption{くぎの中の磁区（模式図。矢印は、それぞれの磁区の小さな磁石が N極を向ける向き、細い線は磁区の境目）。①向きがばらばらで、全体では打ち消し合う。②磁石の磁場と同じ向きの磁区が大きくなり、くぎ全体が磁石になる。}\label{ph:電-7:fig-domain}
\end{bkfigure}

\phboxhead{\bkhatenaref の答え}
\label{ph:電-7:hatena}磁石に近づける前のくぎでは、磁区の向きがばらばらなので、全体では磁石の働きが打ち消し合っています（図\ref{ph:電-7:fig-domain}の①）。磁石を近づけると、磁石の磁場と同じ向きの磁区が、となりの磁区の分まで広がっていきます（磁区の境目が動く）。そうして、くぎ全体で小さな磁石の向きがそろい、くぎそのものが磁石になります（②）。これを\textbf{磁化}といいます。\bkidx{じか}{磁化}磁石のN極にくっついたくぎの頭は S極に、先は N極になるので、くぎの先は次のくぎやクリップを引きつけます。

磁石から離すと、多くの磁区は、もとのばらばらの向きにもどろうとします。けれども、磁区の境目は、鉄の中の小さなきずや不純物に引っかかって、全部はもどれません。だから、くぎには少し磁石の働きが残ります（やってみようの「次の日」）。この「残る」性質は⑧で、磁区のしくみそのものは⑩で、もう一度見ます。アルミや銅には、原子の小さな磁石の向きをそろえる強い働きがなく、磁区ができません。だから、磁石につかないのです（④）。

\begin{phmijika}[キャッシュカードの磁気ストライプ]
キャッシュカードやきっぷの裏の黒い帯（磁気ストライプ）には、とても小さな磁石の粒がぬってあります。カードを読む機械は、帯の場所ごとに、粒の磁石の向きを N・S のどちら向きにそろえて書きこみ、読むときはその向きを調べます。強い磁石を近づけると、書きこんだ向きがそろえ直されて、情報が消えてしまうことがあります。
\end{phmijika}

\Needspace{12\baselineskip}\phstage{te}{磁化と磁気回路を数で扱う}{3}

\bksubsection{3}{磁化 $\vec M$ と磁場の強さ $\vec H$：$\vec B=\mu_0(\vec H+\vec M)$}
\label{ph:電-7:sec3}
1 m あたり $n$ 回巻いた長いソレノイドに電流 $I$ を流すと、中の磁場は $B=\mu_0nI$ でした（電-6 ⑥）。この中に鉄の棒を入れると、磁場は何百倍・何千倍にも強くなります。鉄が磁化して、それ自身が磁場をつくるからです。

\begin{bkfigure}
\begin{tikzpicture}[x=1mm, y=1mm]
  % 左：断面の小さな電流
  \draw[bkInk, line width=0.8pt, fill=bkGray!15] (0,0) rectangle (36,24);
  \foreach \xx in {6,18,30} {\foreach \yy in {6,18} {\draw[bkOrange!80!black, line width=0.8pt, -{Stealth[length=1.5mm]}] (\xx+4,\yy) arc (0:330:4);}}
  \node[sxh lab, anchor=north, text width=40mm, align=center] at (18,-1.5) {鉄の断面（1つ1つの円が\\小さな磁石＝小さな電流）};
  \draw[bkInk, line width=0.9pt, -{Stealth[length=2.4mm]}] (40,12) -- (48,12);
  % 右：表面の電流
  \begin{scope}[shift={(52,0)}]
  \draw[bkInk, line width=0.8pt, fill=bkGray!15] (0,0) rectangle (36,24);
  \draw[bkOrange!85!black, line width=1.6pt, -{Stealth[length=2.6mm]}] (36,0.8) -- (36,23.2);
  \draw[bkOrange!85!black, line width=1.6pt, -{Stealth[length=2.6mm]}] (35.2,24) -- (0.8,24);
  \draw[bkOrange!85!black, line width=1.6pt, -{Stealth[length=2.6mm]}] (0,23.2) -- (0,0.8);
  \draw[bkOrange!85!black, line width=1.6pt, -{Stealth[length=2.6mm]}] (0.8,0) -- (35.2,0);
  \node[sxh lab, text width=30mm, align=center] at (18,12) {中では\\打ち消し合う};
  \node[sxh lab, anchor=north, text width=40mm, align=center] at (18,-1.5) {周りにだけ電流が残る\\（ソレノイドと同じ形）};
  \end{scope}
  \node[anchor=west, font=\small, text width=30mm, align=left] at (94,12) {%
  となりどうしの小さな電流は、くっついた所で向きが反対で打ち消し合う。};
\end{tikzpicture}
\bkfigcaption{磁化した鉄の棒の断面（棒の軸は紙に垂直）。小さな磁石（小さな円い電流）がそろうと、表面を回る電流が残り、ソレノイドの巻いた導線と同じ働きをする。}\label{ph:電-7:fig-surface}
\end{bkfigure}

\noindent 鉄の中の小さな磁石を、電-6 ④の小さな円い電流と考えます。向きがそろうと、図\ref{ph:電-7:fig-surface}のように、となりどうしの電流は打ち消し合い、棒の表面を回る電流だけが残ります。棒の長さ 1 m あたりの、この表面の電流を $M$ と書き、\textbf{磁化}の大きさといいます。単位は、巻数×電流÷長さと同じ A/m です（1 m$^3$ あたりの小さな磁石の強さ $NIS$ の合計と同じ量）。磁化した鉄は、「1 m あたり $M$ アンペアの電流を巻いたソレノイド」と同じ磁場をつくるのです。

導線の電流がつくる分 $nI$ と、鉄の表面の電流がつくる分 $M$ を合わせると、ソレノイドの中の磁場は
\pheq{B=\mu_0(nI+M)}{T ＝ N/A$^2$ × A/m}
です。ここで、導線の電流だけで決まる部分 $nI$ を
\[H=nI\]
と書き、\textbf{磁場の強さ}といいます。\bkidx{じばのつよさ}{磁場の強さ（$H$）}単位は A/m です。電-6 ③で見た高校の教科書の $H$ と同じものです。ベクトルで書くと、
\pheq{\vec B=\mu_0(\vec H+\vec M)}{T ＝ N/A$^2$ × A/m}
\label{ph:電-7:BHM}です。$\vec H$ は「導線を流れる電流だけで決まる部分」、$\vec M$ は「物が磁化してつくる部分」で、力を決めるのは、その和の $\vec B$ です。

\phboxhead{$H$ は、導線の電流だけで決まる}
ソレノイドでは $H=nI$ でした。もっと一般に、閉じた道1周について、$H$ の道にそう部分×長さをたすと、道が囲む\textbf{導線の電流}になります（鉄の表面の電流は数えない）。電-6 ⑥のアンペールの法則の $\vec H$ 版です。
\pheq{\oint_C\vec H\cdot\dd\vec r=I_{\text{中}}}{A ＝ A/m × m}
たとえば、ドーナツの形の鉄心（長さ $l$）に $N$ 回巻いて電流 $I$ を流すと、鉄心にそって1周する道で $Hl=NI$ です。$NI$（アンペアターン）が、$H$ を決めます。

\phboxhead{比透磁率と $B=\mu H$}
多くの物では、磁化 $M$ は、その場所の $H$ に比例します。比例の定数を $\chi$（カイ）と書き、\textbf{磁化率}といいます（単位なし）。
\[M=\chi H\qquad\text{なので}\qquad B=\mu_0(1+\chi)H\]
$1+\chi$ を\textbf{比透磁率} $\mu_{\mathrm r}$、$\mu=\mu_{\mathrm r}\mu_0$ をその物の\textbf{透磁率}といいます。\bkidx{ひとうじりつ}{比透磁率}\bkidx{とうじりつ}{透磁率}
\pheq{B=\mu H=\mu_{\mathrm r}\mu_0H}{T ＝ N/A$^2$ × A/m}
\label{ph:電-7:bmuh}これが、高校の教科書の $B=\mu H$ です。真空では $\mu_{\mathrm r}=1$ で $B=\mu_0H$（電-6 ③）にもどります。空気は $\mu_{\mathrm r}$ が1と10万分の1も違わないので、真空と同じに扱います。鉄心を入れたソレノイドの中は、$B=\mu_{\mathrm r}\mu_0nI$ になり、$\mu_{\mathrm r}$ 倍に強くなります（例題2）。

\begin{phflow}[物の中の磁場]
\phflowrow{$\vec B=\mu_0(\vec H+\vec M)$}{$M$ が $H$ に比例\\（$M=\chi H$）}{$B=\mu H$\\$\mu=(1+\chi)\mu_0$}
\phflowrow{$\oint_C\vec H\cdot\dd\vec r=I_{\text{中}}$}{長いソレノイド}{$H=nI$}
\phflowrow{$B=\mu H$}{鉄心のソレノイド\\$\mu=\mu_{\mathrm r}\mu_0$}{$B=\mu_{\mathrm r}\mu_0nI$}
\end{phflow}

\begin{phyomi}[$\vec B=\mu_0(\vec H+\vec M)$]
「物の中の磁場は、導線の電流がつくる分と、物がつくる分の和」。鉄では $M$ が $H$ の何千倍にもなるので、$B$ はほとんど $\mu_0M$ です（例題1）。鉄が磁場を強くするのは、鉄の中の小さな磁石が向きをそろえて、自分で磁場をつくるからです。
\end{phyomi}

\begin{bkfigure}
\begin{tikzpicture}
\begin{axis}[bkplot, width=100mm, height=50mm, xmin=0, xmax=2200, ymin=0, ymax=2.0, xtick={500,1000,1500,2000}, ytick={0.5,1.0,1.5,2.0}, xlabel={磁場の強さ $H$（A/m）}, ylabel={$B$（T）}, scaled x ticks=false]
  \addplot[bkcurve, domain=0:2200, samples=120] {1.65*tanh(x/400)+4*pi*1e-7*x};
  \addplot[bkOrange, line width=1.4pt, dash pattern=on 4pt off 2pt, domain=0:485] {1.65/400*x};
  \addplot[bkRed!80!black, line width=1.2pt, dash pattern=on 1.2pt off 1.4pt, domain=0:2200] {4*pi*1e-7*x};
  \node[font=\sffamily\footnotesize, anchor=west, inner sep=1pt] at (axis cs:1150,1.80) {鉄（実線）：飽和する};
  \node[font=\sffamily\footnotesize, anchor=west, align=left, inner sep=1pt] at (axis cs:560,0.95) {はじめの傾きの直線（破線）\\$\mu_{\mathrm r}$ 約3300};
  \node[font=\sffamily\footnotesize, anchor=south west, inner sep=1pt] at (axis cs:1500,0.06) {空気（点線）};
\end{axis}
\end{tikzpicture}
\bkfigcaption{鉄の $B$ と $H$ の関係（はじめて磁化するときの曲線。形を見せるための模式図）。はじめは $H$ に比例するが（破線は、はじめの傾きの直線 $B=\mu H$）、$H$ を大きくすると、ほとんどの磁区がそろい終わって、$B$ は増えにくくなる（\textbf{飽和}）。空気の $B=\mu_0H$ は、この目もりでは横軸にほとんど重なる。}\label{ph:電-7:fig-BH0}
\end{bkfigure}

\begin{seikaku}
鉄の $\mu_{\mathrm r}$ は、決まった数ではありません。図\ref{ph:電-7:fig-BH0}のように、$B$ と $H$ は直線の関係ではなく、$H$ を大きくすると飽和します。また、同じ $H$ でも、前にどう磁化したかで $B$ が違います（⑧のヒステリシス）。$B=\mu H$ は、「磁場が弱く、飽和の手前で、ヒステリシスを無視するとき」の近似です。鉄や電磁鋼板の $\mu_{\mathrm r}$ は、材料や磁場の強さで、数百から数千まで大きく変わります。この章の例題の $\mu_{\mathrm r}$ は、計算のために置いた値です。
\end{seikaku}

\bksubsection{4}{強磁性・常磁性・反磁性}
\label{ph:電-7:sec4}
磁石への反応で、物は3つに分けられます。\bkidx{きょうじせい}{強磁性}

\begin{bkcenter}
\begin{bktable}{colspec={Q[l,m,wd=20mm] Q[l,m,wd=32mm] Q[c,m,wd=30mm] Q[l,m,wd=36mm]}, rowsep=1.5pt, row{1}={bg=bkTheme!10, font=\sffamily\small}, cells={font=\small}}
分け方 & 例 & 磁化率 $\chi$ & 磁石への反応 \\
強磁性 & 鉄・ニッケル・コバルト、それらの合金 & 数百〜数千（一定でない） & 強く引きつけられる。磁石にもなる \\
常磁性 & アルミニウム、酸素 & 正でとても小さい（アルミ $+2.2\times10^{-5}$） & ほんの少し引きつけられる \\
反磁性 & 水、銅、炭素、生き物の体 & 負でとても小さい（水 $-9.0\times10^{-6}$、銅 約 $-1.0\times10^{-5}$） & ほんの少し退けられる \\
\end{bktable}
\bktabcaption{磁石への反応による物の分け方（磁化率は、体積あたりの SI の値。20 ℃くらい）}\label{ph:電-7:tab-chi}
\end{bkcenter}

\noindent \textbf{強磁性}の物は、②⑩の磁区をもち、磁化率がとても大きく、磁石から離しても磁化が残ることがあります。\textbf{常磁性}の物は、原子が小さな磁石をもっていますが、となりどうしの向きをそろえる働きがなく、熱の動きでばらばらです。磁場をかけると、ほんの少しだけ磁場の向きにそろいます。\textbf{反磁性}の物は、磁場をかけると、原子の中の電子の動きが少し変わって、磁場と\textbf{反対向き}の磁化ができます（電-8 のレンツの法則と似た働き）。反磁性は、すべての物がもっている性質ですが、とても弱いので、ほかの磁性がない物でしか見えません。\bkidx{はんじせい}{反磁性}

\phboxhead{水も、磁石に退けられる}
磁場の強さが場所で変わる所では、磁化した物は、磁場の強い所へ引かれたり（$\chi>0$）、弱い所へ押し出されたり（$\chi<0$）します。小さな物の 1 m$^3$ あたりにはたらく力は、磁場が上に向かって変わるとき、上向きに
\pheq{f=\frac{\chi}{\mu_0}\,B\,\frac{\dd B}{\dd z}}{N/m$^3$ ＝ $\dfrac{1}{\text{N/A}^2}$ × T × T/m}
です（$|\chi|$ が小さいときの式。導き方は「せいかくに言うと」）。水は $\chi<0$ なので、磁場の弱い向きへ押されます。とても強い磁石の上で、この力が重さ（1 m$^3$ あたり $\rho g$）とつり合うと、水や水でできた生き物が浮かびます。1997年には、\qty{16}{T}の強い電磁石の中で、生きたカエルを浮かせる実験が報告されました（例題10）。

\begin{seikaku}
磁化率 $\chi$ の物を体積 $V$ だけ磁場の中に置くと、物が磁化することで、磁場のエネルギーが $-\dfrac{\chi}{2\mu_0}B^2V$ だけ変わります（$|\chi|\ll1$ のとき。磁場のエネルギー密度 $B^2/(2\mu_0)$ は電-8 ⑧）。力は、エネルギーが下がる向きにはたらき、その大きさは、位置で微分して $F=\dfrac{\chi V}{\mu_0}B\dfrac{\dd B}{\dd z}$ です（力-4 の $F=-\dd U/\dd z$ と同じ考え）。
\end{seikaku}

\bksubsection{5}{磁気回路：電気回路と同じ形で解く}
\label{ph:電-7:sec5}
変圧器やモーターや電磁石では、鉄の芯（\textbf{鉄心}）に導線を巻きます。鉄心は空気より何千倍も磁場を通しやすいので、磁場は鉄心の中に集まり、鉄心にそって輪になって流れます。導線の中を電流が流れるのと似ています。\bkidx{てっしん}{鉄心}

\phboxhead{磁束：鉄心の断面を通る磁場の量}
鉄心の断面積を $S$ とし、断面の中の磁場 $B$ が一様だとすると、断面を通る磁場の量は
\[\Phi=BS\]
です。$\Phi$（ファイ）を\textbf{磁束}といい、単位はウェーバ（Wb。$\qty{1}{Wb}=\qty{1}{T.m^2}$）です。\bkidx{じそく}{磁束}$B$ が「1 m$^2$ あたりの磁束」なので、$B$ を「磁束密度」ともよぶのです。磁束は、電-8 ③で電磁誘導を考えるときの主役になります。$\mathrm{div}\,\vec B=0$（電-6 ⑦）なので、鉄心の中を流れる磁束は、とちゅうで増えたり減ったりしません。導線の中の電流が、とちゅうで増えたり減ったりしない（電-4 ⑥）のと同じです。

\begin{bkfigure}
\begin{tikzpicture}[x=1mm, y=1mm]
  % 磁気回路
  \fill[sxh core] (0,0) rectangle (44,36) (8,8) rectangle (36,28);
  \foreach \yy in {11,14,...,26} {\draw[bkOrange!80!black, line width=1.1pt] (-1.5,\yy) -- (9.5,\yy+1.2);}
  \draw[bkOrange!85!black, line width=1pt] (-1.5,11) -- (-8,11) (-1.5,26) -- (-8,26);
  \node[sxh lab, anchor=east] at (-8.5,18.5) {$I$};
  \draw[sxh I] (-6,12) -- (-6,20);
  \node[sxh lab, anchor=west] at (10,18.5) {$N$ 回};
  \draw[bkBlue!80!black, line width=1.2pt, dash pattern=on 3pt off 1.5pt] (4,4) rectangle (40,32);
  \draw[sxh B, line width=1.2pt] (20,32) -- (24,32);
  \draw[sxh B, line width=1.2pt] (24,4) -- (20,4);
  \node[sxh lab, anchor=south] at (22,36.5) {磁束 $\Phi$（長さ $l$、断面 $S$）};
  \node[sxh lab, anchor=north] at (22,-1.5) {磁気回路};
  % 等号
  \node[font=\sffamily\bfseries\Large, text=bkDeep] at (56,18) {＝};
  \node[font=\sffamily\bfseries\footnotesize, text=bkDeep] at (56,11) {同じ形};
  % 電気回路
  \begin{scope}[shift={(70,0)}]
  \draw[bkInk, line width=0.9pt] (0,4) -- (0,14) (0,22) -- (0,32) -- (13,32) (31,32) -- (40,32) -- (40,4) -- (0,4);
  \draw[bkInk, line width=1.8pt] (-2,14) -- (2,14);
  \draw[bkInk, line width=0.9pt] (-4,22) -- (4,22);
  \node[sxh lab, anchor=east] at (-4.5,18) {$\mathcal E$};
  \draw[bkInk, line width=0.9pt] (13,29) rectangle (31,35);
  \node[sxh lab] at (22,32) {$R$};
  \draw[sxh I] (40,24) -- (40,12);
  \node[sxh lab, anchor=west] at (41,18) {$I$};
  \node[sxh lab, anchor=north] at (20,-1.5) {電気回路};
  \end{scope}
\end{tikzpicture}
\bkfigcaption{磁気回路（左）と電気回路（右）。巻いた導線のアンペアターン $NI$ が、電池の起電力 $\mathcal E$ にあたり、鉄心を回る磁束 $\Phi$ が、導線を回る電流 $I$ にあたる。}\label{ph:電-7:fig-circuit}
\end{bkfigure}

\phboxhead{磁気回路のオームの法則}
断面積 $S$、長さ $l$、透磁率 $\mu$ のドーナツの形の鉄心に、導線を $N$ 回巻いて、電流 $I$ を流します。磁場が鉄心の外へ漏れないとすると、鉄心にそって1周する道で $Hl=NI$（③）。$B=\mu H$、$\Phi=BS$ なので、
\[\Phi=BS=\mu HS=\mu S\,\frac{NI}{l}\]
です。これを書き直すと、
\pheq{NI=\mathcal R_{\mathrm m}\Phi,\qquad\mathcal R_{\mathrm m}=\frac{l}{\mu S}}{A ＝ 1/H × Wb}
\label{ph:電-7:magcircuit}となります。電気回路のオームの法則 $V=RI$、抵抗 $R=l/(\sigma S)$（電-4 ③。$\sigma$ は導電率）と、まったく同じ形です。$NI$〔A〕を\textbf{起磁力}、$\mathcal R_{\mathrm m}$〔1/H〕を\textbf{磁気抵抗}といいます（H は電-8 で学ぶインダクタンスの単位ヘンリー。1 A/Wb は 1/H と同じ）。\bkidx{きじりょく}{起磁力}\bkidx{じきていこう}{磁気抵抗}

\begin{phyomi}[$NI=\mathcal R_{\mathrm m}\Phi$]
「巻数×電流が、磁束を鉄心に押し流す。流れにくさは、長く・細く・透磁率の小さい道ほど大きい」。電気回路の解き方（直列・並列、キルヒホッフの法則）が、そのまま使えます（対応は表\ref{ph:電-7:tab-souji}）。
\end{phyomi}

\begin{phbaisuru}{電流 $I$}
磁束も2倍になります（飽和するまで）。巻数を2倍にしても同じです。鉄心の長さを2倍にすると、磁気抵抗が2倍になり、磁束は半分になります。
\end{phbaisuru}

\begin{bkcenter}
\begin{bktable}{colspec={Q[l,m,wd=26mm] Q[l,m,wd=44mm] Q[l,m,wd=44mm]}, rowsep=1.5pt, row{1}={bg=bkTheme!10, font=\sffamily\small}, cells={font=\small}}
役わり & 電気回路 & 磁気回路 \\
押すもの & 起電力 $\mathcal E$〔V〕 & 起磁力 $NI$〔A〕 \\
流れるもの & 電流 $I$〔A〕 & 磁束 $\Phi$〔Wb〕 \\
流れにくさ & 抵抗 $R=\dfrac{l}{\sigma S}$〔$\Omega$〕 & 磁気抵抗 $\mathcal R_{\mathrm m}=\dfrac{l}{\mu S}$〔1/H〕 \\
通しやすさ & 導電率 $\sigma$ & 透磁率 $\mu$ \\
法則 & $\mathcal E=RI$、直列は和、並列は逆数の和 & $NI=\mathcal R_{\mathrm m}\Phi$、直列は和、並列は逆数の和 \\
分かれ目 & 電流の和は0（電-4 ⑦） & 磁束の和は0（$\mathrm{div}\,\vec B=0$） \\
エネルギー & 抵抗は熱にする & 磁気抵抗は熱にせず、ためる \\
\end{bktable}
\bktabcaption{電気回路と磁気回路の相似。前付の「相似の地図」の表（押すもの・流れるもの）に、磁気回路の列を足したもの。}\label{ph:電-7:tab-souji}
\end{bkcenter}

\begin{phflow}[磁気回路]
\phflowrow{$\oint_C\vec H\cdot\dd\vec r=I_{\text{中}}$}{細い環の鉄心\\漏れなし}{$Hl=NI$}
\phflowrow{$Hl=NI$、$B=\mu H$}{$\Phi=BS$\\断面で一様}{$NI=\mathcal R_{\mathrm m}\Phi$\\$\mathcal R_{\mathrm m}=\dfrac{l}{\mu S}$}
\phflowrow{$NI=\mathcal R_{\mathrm m}\Phi$}{すき間がある\\（直列）}{$NI=(\mathcal R_{\text{鉄}}+\mathcal R_{\text{すき間}})\Phi$}
\end{phflow}

\begin{seikaku}
相似の表の最後の行のように、磁気抵抗は電気抵抗と違って、熱を出しません。鉄心の磁場には、エネルギー $\tfrac12\mathcal R_{\mathrm m}\Phi^2$ がたまります（電池の電流が抵抗で熱になるのとは違い、電流を切れば返ってくる）。エネルギーまで考えると、磁束はコンデンサーにたまる電荷に、磁気抵抗は静電容量の逆数にあたる、と見るほうが正確です。けれども、「どれだけの磁束が流れるか」を求めるだけなら、表のように電気抵抗の回路と同じに解いてかまいません。
\end{seikaku}

\bksubsection{6}{すき間のある磁気回路：電磁石とリレー}
\label{ph:電-7:sec6}
電磁石や、リレー（電磁石で接点を開け閉めするスイッチ）、モーターでは、鉄心に\textbf{すき間}（ギャップ）があります。鉄が引きつけられて動くためには、すき間が要るからです。すき間の空気（$\mu_{\mathrm r}=1$）は、鉄の何千倍も磁場を通しにくいので、すき間が短くても、磁気抵抗の大部分はすき間がしめます。\bkidx{すきま}{すき間（ギャップ）}

\begin{bkfigure}
\begin{tikzpicture}[x=1mm, y=1mm]
  \fill[sxh core] (0,0) rectangle (40,32) (8,8) rectangle (32,24);
  \fill[white] (31.6,15) rectangle (40.4,17.6);
  \draw[bkInk, line width=0.6pt] (32,15) -- (40,15) (32,17.6) -- (40,17.6);
  \foreach \yy in {9,12,...,22} {\draw[bkOrange!80!black, line width=1.1pt] (-1.5,\yy) -- (9.5,\yy+1.2);}
  \node[sxh lab, anchor=east] at (-2,16) {$N$ 回};
  \draw[bkInk, line width=0.5pt] (40.5,15) -- (43,15) (40.5,17.6) -- (43,17.6);
  \node[sxh lab, anchor=west] at (43.2,16.3) {$l_{\mathrm g}$};
  \draw[bkBlue!80!black, line width=1.2pt, dash pattern=on 3pt off 1.5pt] (4,4) rectangle (36.5,28);
  \node[sxh lab, anchor=north] at (20,-1.5) {すき間のある鉄心};
  \node[anchor=west, font=\small, text width=74mm, align=left] at (50,16) {%
  鉄の部分（長さ $l$）とすき間（長さ $l_{\mathrm g}$）は、磁束にとって\textbf{直列}につながった2つの磁気抵抗。\par\vspace{2pt}
  $\mathcal R_{\text{鉄}}=\dfrac{l}{\mu_{\mathrm r}\mu_0S}$、$\mathcal R_{\text{すき間}}=\dfrac{l_{\mathrm g}}{\mu_0S}$\par\vspace{2pt}
  $\mu_{\mathrm r}=2000$ なら、すき間\qty{1}{mm}の磁気抵抗は、鉄\qty{2}{m}分と同じ。};
\end{tikzpicture}
\bkfigcaption{すき間のある磁気回路。すき間の磁束は、鉄心と同じ（すき間のはしで少しふくらむ漏れは無視する）。}\label{ph:電-7:fig-gap}
\end{bkfigure}

\noindent 磁束は直列の磁気抵抗を順に通るので、
\pheq{NI=(\mathcal R_{\text{鉄}}+\mathcal R_{\text{すき間}})\,\Phi}{A ＝ 1/H × Wb}
\label{ph:電-7:gap}です。すき間の磁気抵抗が鉄より十分大きいとき（$l_{\mathrm g}\gg l/\mu_{\mathrm r}$）は、鉄の分を無視して、$NI\approx\Phi l_{\mathrm g}/(\mu_0S)$、つまりすき間の磁場は $B\approx\mu_0NI/l_{\mathrm g}$ です。

\begin{phbaisuru}{すき間 $l_{\mathrm g}$}
すき間の磁気抵抗が大部分をしめるときは、磁束はほぼ半分になります。下のグラフは、例題3・4の鉄心（$\mu_{\mathrm r}=2000$、長さ\qty{0.2}{m}、200回巻き、\qty{0.5}{A}）の磁束密度です。すき間が0のときは\qty{1.26}{T}、\qty{0.1}{mm}あくだけで半分になります。
\par\smallskip\noindent\begin{tikzpicture}
\begin{axis}[bkplot, width=96mm, height=44mm, xmin=0, xmax=2.1, ymin=0, ymax=1.4, xtick={0.5,1,1.5,2}, ytick={0.5,1.0}, xlabel={すき間 $l_{\mathrm g}$（mm）}, ylabel={$B$（T）}]
  \addplot[bkcurve, domain=0:2.05, samples=120] {4*pi*1e-7*100/(0.2/2000+x*1e-3)};
  \fill[bkRed] (axis cs:1,0.1142) circle (1.8pt) (axis cs:2,0.0598) circle (1.8pt) (axis cs:0,1.2566) circle (1.8pt);
  \node[font=\sffamily\footnotesize, anchor=south west, text=bkRed!75!black] at (axis cs:1.05,0.2) {\qty{1}{mm}で\qty{0.11}{T}、\qty{2}{mm}で\qty{0.060}{T}};
\end{axis}
\end{tikzpicture}
\end{phbaisuru}

\begin{phmijika}[電磁石のクレーンと、ドアの電磁ロック]
鉄くずをつり上げる電磁石のクレーンは、電磁石と鉄くずの間のすき間を小さくするほど、磁束が増えて強く引きつけます。ぴったりくっつくと、すき間の磁気抵抗がなくなって磁束がいちばん大きくなります。電気を切ると、鉄心の磁化がほとんど消えて、鉄くずを放します（鉄心には、磁化が残りにくい鉄を使う：⑧）。
\end{phmijika}

\begin{phmodel}[すき間のある磁気回路]
\phignore{鉄心の外へ漏れる磁束（漏れ磁束）。すき間のはしで磁力線がふくらむこと。鉄心の飽和とヒステリシス。}
\phassume{磁束は鉄心の断面に一様に流れる。鉄の $\mu_{\mathrm r}$ は一定。鉄心の長さは、断面の中心を通る道の長さで代表させる。}
\phlimit{すき間が鉄心の太さと同じくらい長いとき（漏れが大きい）。鉄の磁束密度が\qty{1.5}{T}をこえて飽和に近いとき。}
\end{phmodel}

\phkokomade

\Needspace{10\baselineskip}\phstage{shiki}{境目・ヒステリシス・磁石を式で}{5}

\bksubsection{7}{境目での磁場：$B$ の法線の部分と $H$ の接線の部分}
\label{ph:電-7:sec7}
鉄と空気の境目で、磁場はどうつながっているのでしょう。電-3 ⑨で電場について考えたのと同じように、境目をはさむ「とても薄い袋」と「とても細い道」を考えます（8-11 の袋と道を、境目にはりつくまで薄く・細くしたもの）。

\begin{bkfigure}
\begin{tikzpicture}[x=1mm, y=1mm]
  \fill[bkGray!25] (0,0) rectangle (60,14);
  \fill[bkBlue!6] (0,14) rectangle (60,32);
  \draw[bkInk, line width=0.8pt] (0,14) -- (60,14);
  \node[sxh lab, anchor=west] at (1,3) {鉄（$\mu_{\mathrm r}$ が大きい）};
  \node[sxh lab, anchor=west] at (1,29) {空気};
  % 薄い袋
  \draw[bkRed!80!black, line width=0.9pt] (8,11.5) rectangle (20,16.5);
  \node[sxh lab, anchor=south] at (14,17) {薄い袋};
  % 細い道
  \draw[bkGreen!50!black, line width=0.9pt, -{Stealth[length=1.6mm]}] (30,12) -- (44,12) -- (44,16) -- (30,16) -- cycle;
  \node[sxh lab, anchor=south] at (37,17) {細い道};
  % 磁力線
  \draw[sxh B, line width=1.2pt] (50,2) -- (56,13.9);
  \draw[sxh B, line width=1.2pt] (56,14.1) -- (56.6,31);
  \draw[bkGray, line width=0.4pt, dashed] (56,4) -- (56,30);
  \node[sxh lab, anchor=north west] at (56.5,9) {$\theta_1$};
  \node[sxh lab, anchor=south west] at (57,20) {$\theta_2$};
  \node[anchor=west, font=\small, text width=62mm, align=left] at (64,16) {%
  袋：入る磁束＝出る磁束（$\mathrm{div}\,\vec B=0$）なので、\textbf{$B$ の境目に垂直な部分}は、両側で等しい。\par\vspace{2pt}
  道：境目に電流が流れていなければ、1周の $H$ の和は0なので、\textbf{$H$ の境目にそう部分}は、両側で等しい。};
\end{tikzpicture}
\bkfigcaption{境目をはさむ薄い袋と細い道。$\theta_1$・$\theta_2$ は、磁力線が境目の法線（垂直な向き）となす角。}\label{ph:電-7:fig-boundary}
\end{bkfigure}

\noindent 2つの条件を式で書くと、物1（鉄）と物2（空気）で、
\[B_1\cos\theta_1=B_2\cos\theta_2,\qquad H_1\sin\theta_1=H_2\sin\theta_2\]
です。$B=\mu H$ を使って2つめを1つめでわると、$\tan\theta_1/\mu_1=\tan\theta_2/\mu_2$、つまり
\pheq{\frac{\tan\theta_2}{\tan\theta_1}=\frac{\mu_2}{\mu_1}}{単位なし}
\label{ph:電-7:refraction}です。磁力線は、境目で折れ曲がります（電-3 ⑨の電気力線の屈折と同じ形）。鉄（$\mu_1$ が大きい）から空気（$\mu_2=\mu_0$）へ出るとき、$\tan\theta_2$ は $\tan\theta_1$ の $\mu_{\mathrm r}$ 分の1になるので、鉄の中で磁力線がどんな向きでも、空気へ出ると\textbf{ほとんど境目に垂直}になります（例題8）。

\begin{phyomi}[$\dfrac{\tan\theta_2}{\tan\theta_1}=\dfrac{\mu_2}{\mu_1}$]
「透磁率の大きい物の中では、磁力線は境目にそって寝ていてもよいが、透磁率の小さい側へ出ると、境目に垂直に立つ」。電磁石や磁石の磁極の面から、磁力線がほぼ垂直に出ているのは、このためです。
\end{phyomi}

\bksubsection{8}{ヒステリシス：$BH$ 曲線の面積が1周の損失}
\label{ph:電-7:sec8}
鉄心に巻いたコイルに交流の電流を流すと、$H$ は正と負を行ったり来たりします。このときの $B$ の変わり方を、横軸 $H$、たて軸 $B$ のグラフにかくと、行きと帰りで違う道を通る輪（ループ）になります。これを\textbf{ヒステリシス}（「おくれ」という意味のギリシャ語から）といいます。\bkidx{ひすてりしす}{ヒステリシス}

\begin{bkfigure}
\begin{tikzpicture}
\begin{axis}[bkplot, width=110mm, height=58mm, xmin=-260, xmax=260, ymin=-1.8, ymax=1.8, xtick={-200,-100,0,100,200}, ytick={-1.5,-1,-0.5,0.5,1,1.5}, xlabel={$H$（A/m）}, ylabel={$B$（T）}, axis lines=left]
  \draw[bkGray, line width=0.6pt] (axis cs:-260,0) -- (axis cs:260,0) (axis cs:0,-1.8) -- (axis cs:0,1.8);
  \fill[bkOrange!15] (axis cs:-250,-1.5499) -- (axis cs:-240,-1.5498) -- (axis cs:-230,-1.5496) -- (axis cs:-220,-1.5494) -- (axis cs:-210,-1.5490) -- (axis cs:-200,-1.5483) -- (axis cs:-190,-1.5472) -- (axis cs:-180,-1.5453) -- (axis cs:-170,-1.5423) -- (axis cs:-160,-1.5374) -- (axis cs:-150,-1.5293) -- (axis cs:-140,-1.5159) -- (axis cs:-130,-1.4942) -- (axis cs:-120,-1.4591) -- (axis cs:-110,-1.4030) -- (axis cs:-100,-1.3148) -- (axis cs:-90,-1.1805) -- (axis cs:-80,-0.9845) -- (axis cs:-70,-0.7163) -- (axis cs:-60,-0.3796) -- (axis cs:-50,0.0000) -- (axis cs:-40,0.3796) -- (axis cs:-30,0.7163) -- (axis cs:-20,0.9845) -- (axis cs:-10,1.1805) -- (axis cs:0,1.3148) -- (axis cs:10,1.4030) -- (axis cs:20,1.4591) -- (axis cs:30,1.4942) -- (axis cs:40,1.5159) -- (axis cs:50,1.5293) -- (axis cs:60,1.5374) -- (axis cs:70,1.5423) -- (axis cs:80,1.5453) -- (axis cs:90,1.5472) -- (axis cs:100,1.5483) -- (axis cs:110,1.5490) -- (axis cs:120,1.5494) -- (axis cs:130,1.5496) -- (axis cs:140,1.5498) -- (axis cs:150,1.5499) -- (axis cs:160,1.5499) -- (axis cs:170,1.5499) -- (axis cs:180,1.5500) -- (axis cs:190,1.5500) -- (axis cs:200,1.5500) -- (axis cs:210,1.5500) -- (axis cs:220,1.5500) -- (axis cs:230,1.5500) -- (axis cs:240,1.5500) -- (axis cs:250,1.5500) -- (axis cs:250,1.5499) -- (axis cs:240,1.5498) -- (axis cs:230,1.5496) -- (axis cs:220,1.5494) -- (axis cs:210,1.5490) -- (axis cs:200,1.5483) -- (axis cs:190,1.5472) -- (axis cs:180,1.5453) -- (axis cs:170,1.5423) -- (axis cs:160,1.5374) -- (axis cs:150,1.5293) -- (axis cs:140,1.5159) -- (axis cs:130,1.4942) -- (axis cs:120,1.4591) -- (axis cs:110,1.4030) -- (axis cs:100,1.3148) -- (axis cs:90,1.1805) -- (axis cs:80,0.9845) -- (axis cs:70,0.7163) -- (axis cs:60,0.3796) -- (axis cs:50,0.0000) -- (axis cs:40,-0.3796) -- (axis cs:30,-0.7163) -- (axis cs:20,-0.9845) -- (axis cs:10,-1.1805) -- (axis cs:0,-1.3148) -- (axis cs:-10,-1.4030) -- (axis cs:-20,-1.4591) -- (axis cs:-30,-1.4942) -- (axis cs:-40,-1.5159) -- (axis cs:-50,-1.5293) -- (axis cs:-60,-1.5374) -- (axis cs:-70,-1.5423) -- (axis cs:-80,-1.5453) -- (axis cs:-90,-1.5472) -- (axis cs:-100,-1.5483) -- (axis cs:-110,-1.5490) -- (axis cs:-120,-1.5494) -- (axis cs:-130,-1.5496) -- (axis cs:-140,-1.5498) -- (axis cs:-150,-1.5499) -- (axis cs:-160,-1.5499) -- (axis cs:-170,-1.5499) -- (axis cs:-180,-1.5500) -- (axis cs:-190,-1.5500) -- (axis cs:-200,-1.5500) -- (axis cs:-210,-1.5500) -- (axis cs:-220,-1.5500) -- (axis cs:-230,-1.5500) -- (axis cs:-240,-1.5500) -- (axis cs:-250,-1.5500) -- cycle;
  \addplot[bkcurve, domain=-250:250, samples=120] {1.55*tanh((x+50)/40)};
  \addplot[bkcurve, domain=-250:250, samples=120] {1.55*tanh((x-50)/40)};
  \draw[bkBlue!80!black, line width=1pt, -{Stealth[length=2mm]}] (axis cs:150,1.548) -- (axis cs:100,1.548);
  \draw[bkBlue!80!black, line width=1pt, -{Stealth[length=2mm]}] (axis cs:-150,-1.548) -- (axis cs:-100,-1.548);
  \fill[bkRed] (axis cs:0,1.315) circle (1.8pt) (axis cs:-50,0) circle (1.8pt) (axis cs:250,1.55) circle (1.8pt);
  \node[font=\sffamily\footnotesize, anchor=south east, fill=white, inner sep=1pt] at (axis cs:-20,1.36) {残留磁束密度 $B_{\mathrm r}$};
  \node[font=\sffamily\footnotesize, anchor=east, fill=white, inner sep=1pt] at (axis cs:-58,0.25) {保磁力 $-H_{\mathrm c}$};
  \node[font=\sffamily\footnotesize, anchor=north, fill=white, inner sep=1pt] at (axis cs:200,1.45) {$B_{\max}$};
  \node[font=\sffamily\footnotesize, anchor=west, fill=white, inner sep=1pt] at (axis cs:60,-0.9) {色の面積＝1周の損失（1 m$^3$ あたり）};
\end{axis}
\end{tikzpicture}
\bkfigcaption{ヒステリシスのループ（形を見せるための模式図。$B_{\max}=1.55$\,T、$H_{\mathrm c}=\qty{50}{A/m}$）。上の枝は $H$ を減らすとき、下の枝は増やすとき。矢印の向きに回る。}\label{ph:電-7:fig-loop}
\end{bkfigure}

\noindent $H$ を0にもどしても、$B$ は0にもどらず、\textbf{残留磁束密度} $B_{\mathrm r}$ が残ります（くぎに磁石の働きが残ったのと同じ：②）。$B$ を0にするには、反対向きに\textbf{保磁力} $H_{\mathrm c}$ だけの $H$ をかけなければなりません。\bkidx{ざんりゅうじそくみつど}{残留磁束密度}\bkidx{ほじりょく}{保磁力}磁区の境目が、きずに引っかかりながら、行きと帰りで違う動き方をするためです。

\phboxhead{ループの面積が、1周で熱になるエネルギー}
コイル（$N$ 回）に電流 $i$ を流して、鉄心（長さ $l$、断面 $S$）の $B$ を変えると、電-8 ③のファラデーの法則で、コイルに電圧 $NS\,\dd B/\dd t$ が生まれます。電源は、これにさからって電流を流すので、1秒に $NS\dfrac{\dd B}{\dd t}\times i$ の仕事をします。$Hl=Ni$（③）を使うと、これは $(lS)\,H\dfrac{\dd B}{\dd t}$ です。$lS$ は鉄心の体積なので、1 m$^3$ あたり、$B$ が $\dd B$ だけ変わる間に、電源は $H\,\dd B$ の仕事をします。ループを1周すると、
\pheq{w=\oint H\,\dd B=\text{ループの面積}}{J/m$^3$ ＝ A/m × T}
\label{ph:電-7:hysteresis}です（5-12の「グラフの面積が仕事」と同じ考え）。行きと帰りで面積が打ち消し合わない分が、1周ごとに鉄心の中で熱になります。これを\textbf{ヒステリシス損}といいます。周波数 $f$ の交流なら、1秒に $f$ 周するので、1 m$^3$ あたりの損失は $wf$ で、\textbf{周波数に比例}します。\bkidx{ひすてりしすそん}{ヒステリシス損}

\phboxhead{表計算でループの面積を求める}
図\ref{ph:電-7:fig-loop}のループの上の枝と下の枝の $B$ を、$H$ を\qty{50}{A/m}ごとに表にし、2つの差（ループのはば）に $\Delta H=\qty{50}{A/m}$ をかけて、台形の和（5-12②）をとります。

\begin{bkcenter}
\begin{bktable}{colspec={Q[c,m,wd=20mm] *{6}{Q[c,m,wd=13mm]}}, rowsep=1pt, row{1}={bg=bkTheme!10, font=\sffamily\footnotesize}, cells={font=\sffamily\footnotesize}}
$H$〔A/m〕 & $-150$ & $-100$ & $-50$ & 0 & 50 & 100 \\
上の枝 $B$〔T〕 & $-1.529$ & $-1.315$ & 0 & 1.315 & 1.529 & 1.548 \\
下の枝 $B$〔T〕 & $-1.550$ & $-1.548$ & $-1.529$ & $-1.315$ & 0 & 1.315 \\
はば〔T〕 & 0.021 & 0.233 & 1.529 & 2.630 & 1.529 & 0.233 \\
\end{bktable}
\bktabcaption{表計算の一部（$H$ が $-250$〜$\qty{250}{A/m}$ のうち、はばが大きい所）。はばの台形の和に\qty{50}{A/m}をかけると、約\qty{310}{J/m^3}。}\label{ph:電-7:tab-sheet}
\end{bkcenter}

\noindent 全部の区間をたすと約\qty{310}{J/m^3}になりました。ループを、はば $2H_{\mathrm c}$、高さ $2B_{\max}$ の長方形と見ると、面積は $4B_{\max}H_{\mathrm c}=4\times1.55\times50=\qty{310}{J/m^3}$ で、ほとんど同じです。角ばったループでは、この「長方形の見積もり」がよく合います（例題7）。

\begin{phbaisuru}{周波数 $f$}
ヒステリシス損（1秒あたり）も2倍になります。1周の損失 $w$ は変わらず、1秒に回る回数が2倍になるからです。電気自動車のモーターのように何百 Hz で磁化する鉄心では、損失が大きくなるので、保磁力の小さい（ループの細い）材料が要ります（章末問題）。
\end{phbaisuru}

\phboxhead{軟らかい磁性材料と、硬い磁性材料}
変圧器やモーターの鉄心には、保磁力が小さく、ループの細い材料（\textbf{軟磁性材料}：ケイ素を混ぜた鉄の板＝\textbf{電磁鋼板}、フェライトなど）を使い、損失を小さくします。電磁鋼板は、約\qty{2}{T}で飽和します。永久磁石には、保磁力が大きく、ループの太い材料（\textbf{硬磁性材料}）を使い、一度磁化したら、ちょっとやそっとの磁場では磁化が消えないようにします（⑨）。

\bksubsection{9}{永久磁石の動作点}
\label{ph:電-7:sec9}
永久磁石は、外から電流を流さなくても、すき間に磁場をつくります。長さ $l_{\mathrm m}$ の磁石と、長さ $l_{\mathrm g}$ のすき間が、断面の同じ鉄の道でつながっているとします（鉄の部分の磁気抵抗は無視）。導線の電流はないので、1周の $H$ の和は0です（③）。
\[H_{\mathrm m}l_{\mathrm m}+H_{\mathrm g}l_{\mathrm g}=0\]
すき間では $H_{\mathrm g}=B/\mu_0$、磁束は磁石とすき間で同じ（断面も同じ）なので、$B$ も同じです。だから、磁石の中の $H_{\mathrm m}$ は
\pheq{H_{\mathrm m}=-\frac{l_{\mathrm g}}{\mu_0l_{\mathrm m}}\,B}{A/m ＝ m÷(N/A$^2$×m) × T}
で、\textbf{負}（磁化と反対向き）です。磁石は、自分のつくる磁場で、少し磁化を弱められる側で働いているのです。

\noindent ネオジム磁石やフェライト磁石では、ループの左上の部分（\textbf{減磁曲線}）が、ほぼ直線 $B=B_{\mathrm r}+\mu_{\mathrm{rec}}\mu_0H$ になります（$\mu_{\mathrm{rec}}$ はリコイル比透磁率で、1に近い）。この直線と上の式の直線の交わる所が、磁石の\textbf{動作点}です。\bkidx{どうさてん}{動作点}2つを解くと、
\pheq{B=\frac{B_{\mathrm r}}{1+\mu_{\mathrm{rec}}\,l_{\mathrm g}/l_{\mathrm m}}}{T}
\label{ph:電-7:opoint}です。すき間が磁石の長さに比べて短いほど、$B$ は $B_{\mathrm r}$ に近づきます。

\begin{bkfigure}
\begin{tikzpicture}
\begin{axis}[bkplot, width=104mm, height=54mm, xmin=0, xmax=1050, ymin=0, ymax=1.5, xtick={250,500,750,1000}, ytick={0.4,0.8,1.2}, xlabel={$|H_{\mathrm m}|$（kA/m）}, ylabel={$B$（T）}, scaled x ticks=false]
  \addplot[bkcurve, domain=0:985] {1.3-1.05*4*pi*1e-4*x};
  \addplot[bkOrange, line width=1.5pt, dash pattern=on 4pt off 2pt, domain=0:290] {0.40-1.1*4*pi*1e-4*x};
  \addplot[bkInk, line width=0.8pt, domain=0:225] {5*4*pi*1e-4*x};
  \fill[bkRed] (axis cs:171,1.074) circle (2pt) (axis cs:52.2,0.328) circle (2pt);
  \node[font=\sffamily\footnotesize, anchor=south west, fill=white, inner sep=1pt] at (axis cs:205,1.12) {ネオジム $1.07$\,T};
  \draw[bkInk!70, line width=0.4pt] (axis cs:56,0.34) -- (axis cs:228,0.40);
  \node[font=\sffamily\footnotesize, anchor=west, fill=white, inner sep=1pt] at (axis cs:230,0.40) {フェライト $0.33$\,T};
  \node[font=\sffamily\footnotesize, anchor=south west, fill=white, inner sep=1pt] at (axis cs:560,0.62) {ネオジムの減磁曲線（実線）};
  \node[font=\sffamily\footnotesize, anchor=west, fill=white, inner sep=1pt] at (axis cs:240,1.38) {← すき間で決まる直線（細い線）};
\end{axis}
\end{tikzpicture}
\bkfigcaption{永久磁石の動作点（例題11 の値）。横軸は、磁石の中の $H$（磁化と反対向き）の大きさ。減磁曲線（実線：ネオジム、破線：フェライト）と、すき間で決まる直線の交わる点で働く。}\label{ph:電-7:fig-opoint}
\end{bkfigure}

\begin{phmijika}[モーターの中の磁石]
電気自動車やエアコンのモーターの回転子には、ネオジム磁石が使われています（電-12）。ネオジム磁石は、残留磁束密度が約1.2〜\qty{1.4}{T}と大きく、保磁力も大きいので、小さくても強い磁場をつくれます。ただし、温度が上がると保磁力が下がり、磁化が弱まりやすくなるので、熱に強くする元素（ジスプロシウムなど）を少し混ぜたり、冷やしたりして使います。フェライト磁石は、残留磁束密度が約\qty{0.4}{T}と小さいのですが、安くてさびないので、スピーカーや冷蔵庫の磁石に使われています。
\end{phmijika}

\begin{motto}[5]
\phboxhead{磁石の良さをくらべる：最大エネルギー積}
すき間にたまる磁場のエネルギーは、エネルギー密度 $B^2/(2\mu_0)$（電-8 ⑧）に、すき間の体積 $Sl_{\mathrm g}$ をかけたものです。上の $H_{\mathrm m}l_{\mathrm m}+H_{\mathrm g}l_{\mathrm g}=0$ と $H_{\mathrm g}=B/\mu_0$ を使うと、
\[\frac{B^2}{2\mu_0}\,Sl_{\mathrm g}=\frac12\,B\,|H_{\mathrm m}|\times Sl_{\mathrm m}\]
となります。右辺の $Sl_{\mathrm m}$ は磁石の体積です。同じすき間に同じ磁場をつくるとき、動作点の $B|H_{\mathrm m}|$ が大きいほど、磁石は小さくてすみます。直線の減磁曲線 $B=B_{\mathrm r}-\mu_{\mathrm{rec}}\mu_0|H_{\mathrm m}|$ の上で $B|H_{\mathrm m}|$ がいちばん大きくなるのは、$B=B_{\mathrm r}/2$ のときで、その値は
\[(BH)_{\max}=\frac{B_{\mathrm r}^{\,2}}{4\mu_{\mathrm{rec}}\mu_0}\]
です。これを\textbf{最大エネルギー積}といい、磁石の性能をくらべる目安に使います。\bkidx{さいだいえねるぎーせき}{最大エネルギー積}例題11 の値で計算すると、ネオジム磁石（$B_{\mathrm r}=\qty{1.3}{T}$、$\mu_{\mathrm{rec}}=1.05$）は約\qty{320}{kJ/m^3}、フェライト磁石（\qty{0.40}{T}、1.1）は約\qty{29}{kJ/m^3}で、約11倍です。同じすき間に同じ磁場をつくるのに、ネオジム磁石は体積が約11分の1ですむのです。実際の磁石では、減磁曲線が途中で折れ曲がることがあり、そのときはこの式より小さくなります。
\end{motto}

\bksubsection{10}{磁区のしくみとキュリー温度}
\label{ph:電-7:sec10}
鉄の原子の小さな磁石が、となりどうし同じ向きにそろおうとする働き（\textbf{交換相互作用}）は、電-6 の電流どうしの磁場の力ではありません。磁石どうしの力としては、ずっと強すぎるのです。これは、電子が「同じ状態に2つ入れない」という量子の決まりと、電子どうしの電気の反発から生まれる働きで、物理シリーズの『原子』（原-7・原-8）で学びます。この章では、その結果だけを使います。\bkidx{こうかんそうごさよう}{交換相互作用}

\phboxhead{なぜ、小部屋に分かれるのか}
鉄の全体が1つの向きにそろうと、鉄の外に大きな磁場ができ、その磁場のエネルギーがたまります。向きの違う磁区に分かれると、外の磁場が打ち消し合ってエネルギーが小さくなります。そのかわり、磁区の境目（\textbf{磁壁}）では向きがだんだん回るので、境目を作るためのエネルギーが要ります。この2つのかね合いで、磁区の大きさが決まります。

\phboxhead{キュリー温度}
温度が上がると、原子の熱の動きが、向きをそろえる働きに勝つようになります。ある温度をこえると、磁区がこわれ、強磁性が消えて常磁性になります。この温度を\textbf{キュリー温度}といいます。鉄では約770 ℃（\qty{1043}{K}）です。\bkidx{きゅりーおんど}{キュリー温度}熱した鉄が磁石につかなくなるのは、このためです（冷えると、また磁石につくようになります）。

\bksubsection{11}{地球の磁場は、何がつくっている？}
\label{ph:電-7:sec11}
地球は大きな磁石です（電-6 ⑫）。でも、地球の中心部は数千 ℃ で、鉄のキュリー温度よりずっと高いので、永久磁石ではありえません。地球の中心には、鉄を主にした固体の\textbf{内核}と、その外側の液体の\textbf{外核}があります。電気をよく通す液体の鉄が、外核の中で対流（熱で上下に回る流れ）し、地球の自転の影響でうずを巻くと、その流れが電流をつくり、電流が磁場をつくり、その磁場の中を流れる鉄に電流が生まれる……という具合に、磁場が自分を保ち続けると考えられています（\textbf{ダイナモ}。発電機のしくみ：電-8）。

\begin{phmada}[地球の磁場の N と S は、いつ入れかわる？]
岩石は、できるときに、その時代の地球の磁場の向きに磁化して、それを記録します。世界中の岩石を調べると、地球の磁場の N と S は、過去に何度も入れかわってきたことがわかっています（\textbf{地磁気の逆転}）。いちばん新しい逆転は約77万年前で、その記録がよく残っている千葉県市原市の地層（千葉セクション）にちなんで、2020年に、約77万年前から約13万年前までの時代が「チバニアン」と名付けられました。けれども、逆転がなぜ起きるのか、次にいつ起きるのかは、予測できません。外核の液体の鉄の流れをコンピュータで計算する研究では、逆転のようなことが起きる例も得られていますが、本物の地球の条件で計算することは、まだできていないのです。
\end{phmada}

\Needspace{8\baselineskip}
\begin{phsaki}{磁性の物理と、磁場の数値解析}
磁区・交換相互作用・スピンは、量子力学を使って理解します（原-7・原-8）。ハードディスクの読み取りや、新しいメモリ（MRAM）では、電子のスピンを使う「スピントロニクス」の技術が使われています。また、モーターや変圧器の設計では、鉄の飽和やヒステリシスを入れて、磁場を有限要素法（力-16 と同じ考え）で計算します。どちらも、磁性物理学や電気機器設計の専門の本で学びます。
\end{phsaki}

\begin{senmonsho}[電気機器・磁性材料の専門書の書き方]
\begin{bkhonbun}
磁気回路において，起磁力 $\mathcal F=NI$〔A〕，磁気抵抗 $R_{\mathrm m}=l/\mu S$〔H$^{-1}$〕とすれば，磁束は $\Phi=\mathcal F/R_{\mathrm m}$〔Wb〕で与えられる．磁気抵抗の逆数をパーミアンスという．磁性体の磁化 $\symbfit M$ を用いれば $\symbfit B=\mu_0(\symbfit H+\symbfit M)$ であり，$\symbfit M$ を磁化の強さ〔T〕と定義して $\symbfit B=\mu_0\symbfit H+\symbfit M$ と書く流儀もある．
\end{bkhonbun}
\bkyaku{$\mathcal F$ は起磁力の記号（この本では $NI$ とそのまま書いた）。$R_{\mathrm m}$ はこの本の $\mathcal R_{\mathrm m}$。パーミアンスは「磁束の通しやすさ」で、電気回路のコンダクタンスにあたる。電気機器の本には、磁化を T の単位で書き、$B=\mu_0H+M$ とする書き方もある（その $M$ は、この本の $\mu_0M$）。どちらの書き方か、本ごとに確かめよう。古い本では $H$ をエルステッド（Oe）、$B$ をガウス（G）で書くことがある（1 Oe は約 \qty{79.6}{A/m}、\qty{1}{T}＝$10^4$\,G）。}
\end{senmonsho}

\begin{taisetsu}
\def\tsl#1{\makebox[7.6\zw][l]{\sffamily\bfseries\small\color{\bkTT}#1}}
\tsl{物の中の磁場}$\vec B=\mu_0(\vec H+\vec M)$、$\oint\vec H\cdot\dd\vec r=I_{\text{中}}$（導線の電流だけ）\\[3pt]
\tsl{透磁率}$M=\chi H$、$B=\mu H$、$\mu=\mu_{\mathrm r}\mu_0$、$\mu_{\mathrm r}=1+\chi$\\[3pt]
\tsl{磁気回路}$NI=\mathcal R_{\mathrm m}\Phi$、$\mathcal R_{\mathrm m}=\dfrac{l}{\mu S}$、$\Phi=BS$。直列は和\\[3pt]
\tsl{すき間}$NI=(\mathcal R_{\text{鉄}}+\mathcal R_{\text{すき間}})\Phi$、すき間が主なら $B\approx\dfrac{\mu_0NI}{l_{\mathrm g}}$\\[3pt]
\tsl{境界条件}$B$ の垂直な部分、$H$ の境目にそう部分が等しい。$\dfrac{\tan\theta_2}{\tan\theta_1}=\dfrac{\mu_2}{\mu_1}$\\[3pt]
\tsl{損失・磁石}$w=\oint H\,\dd B\approx4B_{\max}H_{\mathrm c}$、損失 $wf$。動作点 $B=\dfrac{B_{\mathrm r}}{1+\mu_{\mathrm{rec}}l_{\mathrm g}/l_{\mathrm m}}$
\end{taisetsu}

\begin{kotoba}
\small
\bkword{磁区}{原子の小さな磁石の向きがそろった小部屋。}
\bkword{磁化}{物の中の小さな磁石の向きがそろうこと。その大きさ $M$〔A/m〕。}
\bkword{磁場の強さ}{導線の電流だけで決まる部分 $H$〔A/m〕。$\vec B=\mu_0(\vec H+\vec M)$。}
\bkword{透磁率・比透磁率}{$B=\mu H$ の $\mu$。真空の何倍かが $\mu_{\mathrm r}$。}
\bkword{起磁力・磁気抵抗}{磁気回路の「押すもの」$NI$ と「流れにくさ」$l/(\mu S)$。}
\bkword{ヒステリシス}{$H$ を行き来させると、$B$ が行きと帰りで違う道を通ること。}
\bkword{保磁力・残留磁束密度}{磁化を0にするのに要る $H$。$H=0$ で残る $B$。}
\bkword{キュリー温度}{強磁性が消える温度。鉄は約770 ℃。}
\end{kotoba}

\bkkigocards{\bkkigocard{$\vec H$}{エイチ}{磁場の強さ〔A/m〕。高校の教科書の「磁場」。第2巻のエンタルピー $H$ とは別}{magnetic field strength}{電-7}}
{\bkkigocard{$\vec M$}{エム}{磁化〔A/m〕。電-8 の相互インダクタンス $M$〔H〕とは別の量（使う章が違う）}{magnetization}{電-7}}
\bkkigocards{\bkkigocard{$\mu$・$\mu_{\mathrm r}$}{ミュー・ミューアール}{透磁率と比透磁率。$\mu=\mu_{\mathrm r}\mu_0$（$\mu_0$ は電-6）。第1巻の摩擦係数とは別}{permeability}{電-7}}
{\bkkigocard{$\mathcal R_{\mathrm m}$}{アール・エム}{磁気抵抗〔1/H〕$=l/(\mu S)$。電気抵抗 $R$ と書体で分ける}{reluctance}{電-7}}
\bkkigocards{\bkkigocard{$\Phi$}{ファイ}{磁束〔Wb〕（⑤）。$\Phi=BS$。ファラデーの法則（電-8 ③）の主役。電-2 の $\Phi_E$ と分けるときは $\Phi_B$}{magnetic flux}{電-7}}{}

%% ======================== 例題 ========================
\Needspace{10\baselineskip}\bksection{reidai}{例題でたしかめよう}
\label{ph:電-7:ex1}
\begin{phreidai}[鉄の中の磁化の大きさ]
鉄心の中の磁場が\qty{1.5}{T}です。このときの磁化 $M$ を求めましょう。鉄の比透磁率は5000とします。
\phbamen
鉄心の中の磁場のうち、鉄が磁化してつくる部分の大きさを求める場面です。
\phmushi
鉄の $\mu_{\mathrm r}$ は一定とし、飽和は考えません。
\phzu
$B=\qty{1.5}{T}$、$\mu_{\mathrm r}=5000$、$\mu_0=4\pi\times10^{-7}$\,N/A$^2$。
\phriyu
$\vec B=\mu_0(\vec H+\vec M)$ と $B=\mu_{\mathrm r}\mu_0H$（③）から、$H$ と $M$ を分けて求められる。
\phtateru
\[H=\frac{B}{\mu_{\mathrm r}\mu_0},\qquad M=\frac{B}{\mu_0}-H\]
\phtoku
\[H=\frac{1.5}{5000\times4\pi\times10^{-7}}\approx\qty{239}{A/m},\qquad\frac{B}{\mu_0}=\frac{1.5}{4\pi\times10^{-7}}\approx1.19\times10^6\ \text{A/m}\]
\[M=1.19\times10^6-239\approx1.2\times10^6\ \text{A/m}\]
\phtashikame
単位：T÷(N/A$^2$)＝N/(A·m)×A$^2$/N＝A/m。\\
大きさ：$H$ は $M$ の約5000分の1で、$B$ のほとんどは鉄の磁化がつくっています。$M=1.2\times10^6$\,A/m は、1 m あたり120万アンペアの電流を巻いたソレノイドと同じです。導線でこれを作るのは、とてもたいへんです。\\
極端な場合：鉄がなければ（$\mu_{\mathrm r}=1$）、$M=0$。
\kotae 約 $1.2\times10^6$\,A/m
\end{phreidai}
\begin{jibun}
電磁鋼板の中の磁場が\qty{1.8}{T}のとき、磁化はいくらですか（$H$ は小さいので無視する）。
\bkjans{$M\approx B/\mu_0=1.8\div(4\pi\times10^{-7})$ で、約 $1.4\times10^6$\,A/m}
\end{jibun}

\label{ph:電-7:ex2}
\begin{phreidai}[鉄心を入れたソレノイド]
1 m あたり1000回巻いた長いソレノイドに\qty{0.10}{A}を流します。（1）中が空気のとき、（2）比透磁率1000の鉄心を入れたときの、中の磁場を求めましょう。
\phbamen
鉄心が磁場を強くするようすを、数で確かめる場面です。
\phmushi
ソレノイドは十分長く、鉄心はソレノイドの中を全部うめるとします。鉄の $\mu_{\mathrm r}$ は一定とします。
\phzu
$n=1000$ /m、$I=\qty{0.10}{A}$、$\mu_{\mathrm r}=1000$。
\phriyu
$H=nI$ は導線の電流だけで決まる（③）。$B=\mu_{\mathrm r}\mu_0H$。
\phtateru
\[H=nI,\qquad B=\mu_{\mathrm r}\mu_0H\]
\phtoku
$H=1000\times0.10=\qty{100}{A/m}$。\\
（1）$B=\mu_0H=4\pi\times10^{-7}\times100\approx1.3\times10^{-4}$\,T\\
（2）$B=1000\times1.3\times10^{-4}\approx\qty{0.13}{T}$
\phtashikame
単位：N/A$^2$×A/m＝T。\\
大きさ：（2）は（1）の1000倍。くぎの電磁石が、弱い電流でもクリップを引きつけられるのは、鉄心のおかげです。\\
極端な場合：電流を大きくしていくと、鉄は約\qty{2}{T}で飽和し、それ以上は1000倍にはなりません。
\kotae （1）約 $1.3\times10^{-4}$\,T　（2）約\qty{0.13}{T}
\end{phreidai}
\begin{jibun}
1 m あたり500回巻いたソレノイドに\qty{0.20}{A}を流し、比透磁率2000の鉄心を入れます。中の磁場を求めましょう。
\bkjans{$H=100$\,A/m。$B=2000\times4\pi\times10^{-7}\times100$ で、約\qty{0.25}{T}}
\end{jibun}

\label{ph:電-7:ex3}
\begin{phreidai}[鉄心の磁気回路]
長さ（中心を通る道の長さ）\qty{0.20}{m}、断面積\qty{4.0}{cm^2}、比透磁率2000のドーナツの形の鉄心に、200回巻いて\qty{0.50}{A}を流します。磁気抵抗、磁束、鉄心の中の磁場を求めましょう。答えは3けたで書きます。
\phbamen
変圧器や電磁石の鉄心の磁束を、磁気回路で求める場面です。
\phmushi
漏れ磁束、飽和、ヒステリシスは無視します。
\phzu
図\ref{ph:電-7:fig-circuit}の左。$l=\qty{0.20}{m}$、$S=4.0\times10^{-4}$\,m$^2$、$\mu_{\mathrm r}=2000$、$N=200$、$I=\qty{0.50}{A}$。
\phriyu
磁気回路のオームの法則 $NI=\mathcal R_{\mathrm m}\Phi$、$\mathcal R_{\mathrm m}=l/(\mu S)$（⑤）。
\phtateru
\[\mathcal R_{\mathrm m}=\frac{l}{\mu_{\mathrm r}\mu_0S},\qquad\Phi=\frac{NI}{\mathcal R_{\mathrm m}},\qquad B=\frac{\Phi}{S}\]
\phtoku
\[\mathcal R_{\mathrm m}=\frac{0.20}{2000\times4\pi\times10^{-7}\times4.0\times10^{-4}}\approx1.99\times10^5\ \text{/H}\]
\[\Phi=\frac{200\times0.50}{1.99\times10^5}\approx5.03\times10^{-4}\ \text{Wb},\qquad B=\frac{5.03\times10^{-4}}{4.0\times10^{-4}}\approx\qty{1.26}{T}\]
\phtashikame
単位：A÷(1/H)＝A·H＝Wb（$\qty{1}{H}=\qty{1}{Wb/A}$：電-8）。\\
大きさ：\qty{1.26}{T}は、飽和（約\qty{2}{T}）の手前で、変圧器の鉄心のふつうの値です。$H=NI/l=\qty{500}{A/m}$ からも、$B=2000\times4\pi\times10^{-7}\times500\approx\qty{1.26}{T}$ と同じ答えになります。\\
極端な場合：巻数か電流が0なら、磁束は0。
\kotae 磁気抵抗 約 $1.99\times10^5$\,/H、磁束 約 $5.03\times10^{-4}$\,Wb、約\qty{1.26}{T}
\end{phreidai}
\begin{jibun}
長さ\qty{0.30}{m}、断面積\qty{5.0}{cm^2}、比透磁率1000の鉄心に、100回巻いて\qty{1.0}{A}を流します。磁気抵抗、磁束、磁場を求めましょう。
\bkjans{$\mathcal R_{\mathrm m}\approx4.8\times10^5$\,/H、$\Phi\approx2.1\times10^{-4}$\,Wb、$B\approx\qty{0.42}{T}$}
\end{jibun}

\label{ph:電-7:ex4}
\begin{phreidai}[すき間のある鉄心]
例題3の鉄心に、長さ\qty{1.0}{mm}のすき間を作りました。磁束と、すき間の磁場を求めましょう。全体の磁気抵抗のうち、すき間がしめる割合はいくらですか。
\phbamen
鉄心の小さなすき間が、磁束をどれだけ減らすかを調べる場面です。
\phmushi
すき間のはしでの磁束のふくらみ（漏れ）を無視し、すき間の断面も鉄心と同じとします。鉄の部分の長さは、すき間の分を引いても\qty{0.20}{m}のままとします（\qty{1}{mm}は\qty{0.2}{m}に比べて小さい）。
\phzu
図\ref{ph:電-7:fig-gap}。$l_{\mathrm g}=1.0\times10^{-3}$\,m、$\mathcal R_{\text{鉄}}=1.99\times10^5$\,/H（例題3）。途中の数は3けたで書きます。
\phriyu
鉄とすき間は直列の磁気抵抗（⑥）。
\phtateru
\[\mathcal R_{\text{すき間}}=\frac{l_{\mathrm g}}{\mu_0S},\qquad\Phi=\frac{NI}{\mathcal R_{\text{鉄}}+\mathcal R_{\text{すき間}}}\]
\phtoku
\[\mathcal R_{\text{すき間}}=\frac{1.0\times10^{-3}}{4\pi\times10^{-7}\times4.0\times10^{-4}}\approx1.99\times10^6\ \text{/H}\]
\[\Phi=\frac{100}{1.99\times10^5+1.99\times10^6}\approx4.57\times10^{-5}\ \text{Wb}\]
\[B\approx\frac{4.57\times10^{-5}}{4.0\times10^{-4}}\approx\qty{0.114}{T}\]
すき間の割合は $1.99\times10^6\div(2.19\times10^6)\approx0.91$ です。
\phtashikame
単位：Wb、T。\\
大きさ：\qty{1}{mm}のすき間で、磁束は約11分の1になりました。すき間の磁気抵抗が、鉄\qty{0.2}{m}の10倍だからです。\\
極端な場合：すき間を0にすると例題3の\qty{1.26}{T}にもどります。
\kotae 磁束 約 $4.6\times10^{-5}$\,Wb、約\qty{0.11}{T}、すき間の割合 約91\,\%
\end{phreidai}
\begin{jibun}
例題3の鉄心のすき間が\qty{0.50}{mm}のとき、磁束、磁場、すき間の磁気抵抗の割合を求めましょう。
\bkjans{$\mathcal R_{\text{すき間}}\approx9.9\times10^5$\,/H。$\Phi\approx8.4\times10^{-5}$\,Wb、$B\approx\qty{0.21}{T}$、割合 約83\,\%}
\end{jibun}

\label{ph:電-7:ex5}
\begin{phreidai}[すき間に決まった磁場をつくる]
例題4の鉄心（すき間\qty{1.0}{mm}）で、すき間の磁場を\qty{0.50}{T}にするには、何アンペアターン要りますか。
\phbamen
電磁石の設計で、要る巻数×電流を見積もる場面です。
\phmushi
例題4と同じ。$\mu_{\mathrm r}=2000$ は、\qty{0.5}{T}でも変わらないとします。
\phzu
$B=\qty{0.50}{T}$、$l_{\mathrm g}=1.0\times10^{-3}$\,m、$l=\qty{0.20}{m}$。
\phriyu
$H$ の1周の和が $NI$（③）。すき間の $H$ は $B/\mu_0$、鉄の $H$ は $B/(\mu_{\mathrm r}\mu_0)$。
\phtateru
\[NI=\frac{B}{\mu_0}\,l_{\mathrm g}+\frac{B}{\mu_{\mathrm r}\mu_0}\,l\]
\phtoku
\[\frac{0.50}{4\pi\times10^{-7}}\times1.0\times10^{-3}\approx\qty{398}{A},\qquad\frac{0.50}{2000\times4\pi\times10^{-7}}\times0.20\approx\qty{40}{A}\]
合わせて約\qty{438}{A}、2けたで約\qty{440}{A}です（たとえば440回巻いて\qty{1.0}{A}）。
\phtashikame
単位：A/m×m＝A。\\
大きさ：すき間の分が9割をしめます。すき間だけで見積もっても\qty{400}{A}で、1割の違いです。\\
極端な場合：$\mu_{\mathrm r}\to\infty$ なら鉄の分は0で、すき間の分だけになります。
\kotae 約\qty{440}{A}（アンペアターン）
\end{phreidai}
\begin{jibun}
同じ鉄心で、すき間を\qty{2.0}{mm}にし、すき間の磁場を\qty{0.40}{T}にするには、何アンペアターン要りますか。
\bkjans{すき間の分は $B l_{\mathrm g}/\mu_0$ で約\qty{637}{A}、鉄の分は約\qty{32}{A}。合わせて約\qty{670}{A}}
\end{jibun}

\label{ph:電-7:ex6}
\begin{phreidai}[リレーの電磁石]
6-4 で使ったリレーのコイル（\qty{12}{V}・\qty{400}{\ohm}）が2000回巻きで、鉄心のすき間が\qty{0.50}{mm}とします。すき間の磁場を見積もりましょう。
\phbamen
数学 6-4 のリレーで、電流がどれだけの磁場をつくるかを見積もる場面です。
\phmushi
鉄の部分の磁気抵抗と漏れ磁束を無視し、すき間だけで見積もります。電流は十分時間がたって一定になったときの値とします（流れ始めのようすは電-9）。
\phzu
$V=\qty{12}{V}$、$R=\qty{400}{\ohm}$、$N=2000$、$l_{\mathrm g}=5.0\times10^{-4}$\,m。
\phriyu
電流はオームの法則 $I=V/R$（電-4 ③）。すき間が主なら $B\approx\mu_0NI/l_{\mathrm g}$（⑥）。
\phtateru
\[I=\frac VR,\qquad B\approx\frac{\mu_0NI}{l_{\mathrm g}}\]
\phtoku
\[I=\frac{12}{400}=\qty{0.030}{A},\qquad NI=2000\times0.030=\qty{60}{A}\]
\[B\approx\frac{4\pi\times10^{-7}\times60}{5.0\times10^{-4}}\approx\qty{0.15}{T}\]
\phtashikame
単位：N/A$^2$×A÷m＝T。\\
大きさ：\qty{30}{mA}の小さな電流でも、2000回巻いて、すき間を小さくすると、冷蔵庫の磁石より強い磁場ができます。\\
極端な場合：すき間を2倍にすると、磁場はほぼ半分になり、接点を引く力は弱まります。
\kotae 約\qty{0.15}{T}
\end{phreidai}
\begin{jibun}
\qty{5.0}{V}・\qty{125}{\ohm}のリレーのコイルが1000回巻きで、すき間が\qty{0.40}{mm}とします。すき間の磁場を見積もりましょう。
\bkjans{$I=\qty{0.040}{A}$、$NI=\qty{40}{A}$。$B\approx4\pi\times10^{-7}\times40\div(4.0\times10^{-4})$ で、約\qty{0.13}{T}}
\end{jibun}

\label{ph:電-7:ex7}
\begin{phreidai}[変圧器の鉄心のヒステリシス損]
変圧器の鉄心が、$B_{\max}=\qty{1.5}{T}$、保磁力 $H_{\mathrm c}=\qty{50}{A/m}$ のループで、\qty{50}{Hz}で磁化されています。ループを長方形とみなして、1周の損失、1 m$^3$ あたり・1 kg あたりの損失（1秒あたり）を求めましょう。鉄心の密度は\qty{7650}{kg/m^3}とします。
\phbamen
変圧器の鉄心が熱をもつ原因の1つ（ヒステリシス損）を見積もる場面です。
\phmushi
渦電流損（電-8 ⑪）は考えません。ループは、はば $2H_{\mathrm c}$、高さ $2B_{\max}$ の長方形とします。
\phzu
図\ref{ph:電-7:fig-loop}。$B_{\max}=1.5$、$H_{\mathrm c}=50$、$f=50$、密度 $7650$。
\phriyu
1周の損失はループの面積（⑧）。1秒に $f$ 周する。
\phtateru
\[w\approx4B_{\max}H_{\mathrm c},\qquad\text{1 m}^3\text{ あたり }wf,\qquad\text{1 kg あたり }\frac{wf}{\text{密度}}\]
\phtoku
\[w=4\times1.5\times50=\qty{300}{J/m^3},\qquad wf=300\times50=1.5\times10^4\ \text{W/m}^3\]
\[\frac{1.5\times10^4}{7650}\approx\qty{2.0}{W/kg}\]
\phtashikame
単位：T×A/m＝N/(A·m)×A/m＝N/m$^2$＝J/m$^3$。\\
大きさ：電磁鋼板の規格（JIS）で、よく使う「35A300」の鋼板は、\qty{1.5}{T}・\qty{50}{Hz}の鉄損（ヒステリシス損と渦電流損の和）が 1 kg あたり\qty{3.00}{W}以下と決められています。見積もりは、その桁に合っています。\\
極端な場合：$H_{\mathrm c}=0$（ループが線）なら損失は0。
\kotae 1周 \qty{300}{J/m^3}、$1.5\times10^4$\,W/m$^3$、約\qty{2.0}{W/kg}
\end{phreidai}
\begin{jibun}
$B_{\max}=\qty{1.2}{T}$、$H_{\mathrm c}=\qty{30}{A/m}$ の鉄心を\qty{60}{Hz}で磁化します。1 kg あたりの損失を求めましょう（密度\qty{7650}{kg/m^3}）。
\bkjans{$w=4\times1.2\times30=\qty{144}{J/m^3}$、$144\times60\div7650$ で、約\qty{1.1}{W/kg}}
\end{jibun}

\label{ph:電-7:ex8}
\begin{phreidai}[鉄から空気へ出る磁力線]
比透磁率1000の鉄の中で、磁力線が境目の法線（垂直な向き）と80°の角をなしています（境目にそって寝ている）。空気へ出たとき、法線と何度の角をなしますか。
\phbamen
鉄の表面から磁力線がどう出るかを、境界条件で調べる場面です。
\phmushi
境目に電流は流れていないとします。鉄の $\mu_{\mathrm r}$ は一定とします。
\phzu
図\ref{ph:電-7:fig-boundary}。$\theta_1=80^\circ$（鉄）、$\mu_1/\mu_2=1000$。
\phriyu
磁力線の屈折 $\tan\theta_2/\tan\theta_1=\mu_2/\mu_1$（⑦）。
\phtateru
\[\tan\theta_2=\frac{\mu_2}{\mu_1}\tan\theta_1\]
\phtoku
\[\tan\theta_2=\frac{\tan80^\circ}{1000}=\frac{5.67}{1000}\approx0.0057,\qquad\theta_2\approx0.32^\circ\]
\phtashikame
単位：角なので単位なし。\\
大きさ：鉄の中ではほとんど境目にそっていたのに、空気へ出ると、ほとんど垂直（法線から0.3°）になりました。\\
極端な場合：$\theta_1=0$（垂直）なら $\theta_2=0$ で、まっすぐ出ます。
\kotae 法線から約0.32°（ほとんど垂直）
\end{phreidai}
\begin{jibun}
比透磁率500のフェライトの中で、磁力線が法線と60°をなしています。空気へ出たときの角を求めましょう。
\bkjans{$\tan\theta_2=\tan60^\circ\div500\approx0.0035$ で、約0.20°}
\end{jibun}

\label{ph:電-7:ex9}
\begin{phreidai}[地球の磁場の $H$]
東京あたりの地球の磁場の強さを、全体で約 $4.6\times10^{-5}$\,T とします（電-6 ⑫）。磁場の強さ $H$ は何 A/m ですか。
\phbamen
空気の中の磁場を、$B$ と $H$ の2つの単位で表す場面です。
\phmushi
空気の比透磁率は1とします。
\phzu
$B=4.6\times10^{-5}$\,T。
\phriyu
空気（真空）では $B=\mu_0H$（③）。
\phtateru
\[H=\frac{B}{\mu_0}\]
\phtoku
\[H=\frac{4.6\times10^{-5}}{4\pi\times10^{-7}}\approx\qty{37}{A/m}\]
\phtashikame
単位：T÷(N/A$^2$)＝A/m。\\
大きさ：1 m あたり37回巻いて\qty{1}{A}流したソレノイドと同じ強さです。\\
極端な場合：空気の中では、$B$ と $H$ は $\mu_0$ 倍の違いだけで、どちらで書いても同じことを表しています。
\kotae 約\qty{37}{A/m}
\end{phreidai}
\begin{jibun}
地球の磁場の水平の部分 $3.0\times10^{-5}$\,T を、$H$ で表しましょう。
\bkjans{$3.0\times10^{-5}\div(4\pi\times10^{-7})$ で、約\qty{24}{A/m}}
\end{jibun}

\label{ph:電-7:ex10}
\begin{phreidai}[水を磁石で浮かせる]
水（磁化率 $\chi=-9.0\times10^{-6}$、密度\qty{1000}{kg/m^3}）の小さな粒を、磁場で浮かせます。磁場が上へ行くほど弱くなる所で、$B\,\dd B/\dd z$ の大きさがいくらあれば浮かびますか。$B=\qty{16}{T}$ の所なら、磁場は 1 m あたり何 T の割合で弱くなっていればよいですか。$g=\qty{9.8}{m/s^2}$ とします。
\phbamen
反磁性の水が、強い磁場の弱い向きに押し出される力で、重さを支える場面です。
\phmushi
粒は小さく、$|\chi|\ll1$ とします。空気の浮力は無視します。
\phzu
$\chi=-9.0\times10^{-6}$、$\rho=\qty{1000}{kg/m^3}$、$g=9.8$。上向きを $z$ の正の向きとし、$\dd B/\dd z<0$。
\phriyu
1 m$^3$ あたりの力 $f=(\chi/\mu_0)B\,\dd B/\dd z$（④）。$\chi<0$、$\dd B/\dd z<0$ なので上向き。これが重さ $\rho g$ とつり合う。
\phtateru
\[\frac{|\chi|}{\mu_0}\,\Bigl|B\frac{\dd B}{\dd z}\Bigr|=\rho g\qquad\Longrightarrow\qquad\Bigl|B\frac{\dd B}{\dd z}\Bigr|=\frac{\mu_0\rho g}{|\chi|}\]
\phtoku
\[\Bigl|B\frac{\dd B}{\dd z}\Bigr|=\frac{4\pi\times10^{-7}\times1000\times9.8}{9.0\times10^{-6}}\approx1.4\times10^3\ \text{T}^2\text{/m}\]
$B=\qty{16}{T}$ なら、$1.4\times10^3\div16\approx\qty{86}{T/m}$ の割合で弱くなっていればよい（\qty{10}{cm}で約\qty{8.6}{T}弱くなる）。
\phtashikame
単位：N/A$^2$×kg/m$^3$×m/s$^2$＝N$^2$/(A$^2$·m$^3$)。T$^2$/m＝(N/(A·m))$^2$/m＝N$^2$/(A$^2$·m$^3$) で合う。\\
大きさ：研究で使われる値（1360〜1400 T$^2$/m）と合います。ふつうの磁石（\qty{0.1}{T}くらい）では、とても届きません。\\
極端な場合：$\chi=0$ なら、どんな磁場でも浮きません。
\kotae 約 $1.4\times10^3$\,T$^2$/m。\qty{16}{T}の所で 1 m あたり約\qty{86}{T}
\end{phreidai}
\begin{jibun}
月の上（$g=\qty{1.62}{m/s^2}$）で水を浮かせるには、$B\,\dd B/\dd z$ の大きさがいくらあればよいですか。
\bkjans{$\mu_0\rho g/|\chi|$ の $g$ を1.62にして、約\qty{230}{T^2/m}（地球の約6分の1）}
\end{jibun}

\label{ph:電-7:ex11}
\begin{phreidai}[ネオジム磁石とフェライト磁石]
長さ\qty{5.0}{mm}の磁石が、鉄の道でつながった長さ\qty{1.0}{mm}のすき間に磁場をつくります（磁石とすき間の断面は同じ）。ネオジム磁石（$B_{\mathrm r}=\qty{1.3}{T}$、$\mu_{\mathrm{rec}}=1.05$）とフェライト磁石（$B_{\mathrm r}=\qty{0.40}{T}$、$\mu_{\mathrm{rec}}=1.1$）とします。すき間の磁場と、すき間の面が引き合う力の比を求めましょう。
\phbamen
同じ形の磁石の材料を変えたとき、すき間の磁場と力がどれだけ違うかを比べる場面です。
\phmushi
鉄の部分の磁気抵抗と漏れ磁束を無視します。減磁曲線は直線とします。
\phzu
図\ref{ph:電-7:fig-opoint}。$l_{\mathrm g}/l_{\mathrm m}=1.0\div5.0=0.20$。
\phriyu
動作点の式（⑨）。すき間の面が引き合う力は、面積あたり $B^2/(2\mu_0)$（電-8 ⑩）で、$B^2$ に比例する。
\phtateru
\[B=\frac{B_{\mathrm r}}{1+\mu_{\mathrm{rec}}l_{\mathrm g}/l_{\mathrm m}},\qquad\frac{F_{\text{ネオジム}}}{F_{\text{フェライト}}}=\Bigl(\frac{B_{\text{ネオジム}}}{B_{\text{フェライト}}}\Bigr)^2\]
\phtoku
\[B_{\text{ネオジム}}=\frac{1.3}{1+1.05\times0.20}\approx\qty{1.07}{T},\qquad B_{\text{フェライト}}=\frac{0.40}{1+1.1\times0.20}\approx\qty{0.33}{T}\]
\[\text{力の比}=\Bigl(\frac{1.07}{0.33}\Bigr)^2\approx11\]
\phtashikame
単位：T÷（単位なし）＝T。\\
大きさ：力は約11倍。$B_{\mathrm r}$ の比の2乗 $(1.3/0.40)^2\approx11$ とほぼ同じです。1 cm$^2$ あたりでは、ネオジムで約\qty{46}{N}、フェライトで約\qty{4.3}{N}になります（すき間の面の磁場の力）。\\
極端な場合：すき間を0にすると $B=B_{\mathrm r}$ です。
\kotae ネオジム 約\qty{1.07}{T}、フェライト 約\qty{0.33}{T}、力は約11倍
\end{phreidai}
\begin{jibun}
例題11のネオジム磁石を長さ\qty{10}{mm}にしました（すき間\qty{1.0}{mm}）。すき間の磁場を求めましょう。
\bkjans{$1.3\div(1+1.05\times0.10)$ で、約\qty{1.2}{T}}
\end{jibun}

\label{ph:電-7:ex12}
\begin{phreidai}[枝分かれする磁気回路（外鉄形の鉄心）]
日の字の形の鉄心の、まん中の脚（長さ\qty{0.10}{m}、断面\qty{4.0}{cm^2}）に100回巻いて\qty{1.0}{A}を流します。磁束は、左右の外側の道（それぞれ長さ\qty{0.30}{m}、断面\qty{2.0}{cm^2}）に分かれて戻ります。比透磁率は2000とします。まん中と外側の磁束と磁場を、電気回路と同じ形で求めましょう。
\phbamen
変圧器によく使われる形の鉄心で、磁束が枝分かれする場面です。
\phmushi
漏れ磁束・飽和を無視します。脚の長さは、断面の中心を通る道の長さで代表させます。
\phzu
\begin{minipage}[t]{0.5\linewidth}
まん中の脚 $\mathcal R_{\mathrm c}$ に起磁力 $NI=\qty{100}{A}$ があり、外側の道 $\mathcal R_{\mathrm o}$ が2本、\textbf{並列}につながっている。電気回路なら、電池と抵抗 $R_{\mathrm c}$ に、2本の抵抗 $R_{\mathrm o}$ の並列をつないだ形。
\end{minipage}\hfill
\begin{minipage}[t]{0.46\linewidth}\centering
\begin{tikzpicture}[x=0.8mm, y=0.8mm, baseline=(cur.north)]
  \coordinate (cur) at (0,30);
  \fill[sxh core] (0,0) rectangle (60,30) (6,6) rectangle (27,24) (33,6) rectangle (54,24);
  \foreach \yy in {8,11,...,20} {\draw[bkOrange!80!black, line width=1pt] (25.5,\yy) -- (34.5,\yy+1.2);}
  \draw[bkBlue!80!black, line width=1pt, -{Stealth[length=1.6mm]}] (30,4) -- (30,26);
  \draw[bkBlue!80!black, line width=0.8pt, -{Stealth[length=1.6mm]}] (24,27) -- (8,27);
  \draw[bkBlue!80!black, line width=0.8pt, -{Stealth[length=1.6mm]}] (36,27) -- (52,27);
  \node[sxh lab, anchor=north] at (30,-1) {まん中に $N$ 回};
\end{tikzpicture}
\end{minipage}
\phriyu
磁気回路でも、直列は和、並列は逆数の和（⑤の相似の表）。分かれ目では、入る磁束＝出る磁束。
\phtateru
\[\mathcal R_{\text{全体}}=\mathcal R_{\mathrm c}+\frac{\mathcal R_{\mathrm o}}{2},\qquad\Phi=\frac{NI}{\mathcal R_{\text{全体}}},\qquad\text{外側は }\frac{\Phi}{2}\text{ ずつ}\]
\phtoku
\[\mathcal R_{\mathrm c}=\frac{0.10}{2000\times4\pi\times10^{-7}\times4.0\times10^{-4}}\approx9.9\times10^4\ \text{/H}\]
\[\mathcal R_{\mathrm o}=\frac{0.30}{2000\times4\pi\times10^{-7}\times2.0\times10^{-4}}\approx6.0\times10^5\ \text{/H}\]
\[\mathcal R_{\text{全体}}\approx9.9\times10^4+3.0\times10^5\approx4.0\times10^5\ \text{/H}\]
\[\Phi=\frac{100}{4.0\times10^5}\approx2.5\times10^{-4}\ \text{Wb}\]
まん中の磁場は $2.5\times10^{-4}\div(4.0\times10^{-4})\approx\qty{0.63}{T}$。外側は1本 $1.3\times10^{-4}$\,Wb で、断面が半分なので、磁場は同じ約\qty{0.63}{T}です。
\phtashikame
単位：Wb・T。\\
大きさ：外側の脚の断面をまん中の半分にしておくと、どこでも同じ磁場になり、鉄をむだなく使えます。本物の変圧器の鉄心も、このように作ります。\\
極端な場合：外側の道を1本にすると、並列でなくなり、磁気抵抗が増えて、磁束は減ります。
\kotae まん中 約 $2.5\times10^{-4}$\,Wb、外側 約 $1.3\times10^{-4}$\,Wb ずつ。どこも約\qty{0.63}{T}
\end{phreidai}
\begin{jibun}
例題12の鉄心の比透磁率が1000だったら、まん中の磁束と磁場はいくらですか。
\bkjans{磁気抵抗がどれも2倍になるので、磁束は半分の約 $1.3\times10^{-4}$\,Wb、磁場は約\qty{0.31}{T}}
\end{jibun}

%% ======================== まちがえやすいところ ========================
\Needspace{11\baselineskip}\bksection{tsumazuki}{まちがえやすいところをチェック}
\begin{tsumazuki}{$B$ と $H$ は、いつも $\mu_0$ 倍？}
\bkng{鉄の中でも $B=\mu_0H$。}
\bkok{鉄の中では $B=\mu_0(H+M)$。$B=\mu_{\mathrm r}\mu_0H$ で、$\mu_{\mathrm r}$ は数百〜数千。}
$B=\mu_0H$ は真空（空気）の中だけです（③）。
\end{tsumazuki}

\begin{tsumazuki}{磁気抵抗も熱を出す？}
\bkng{磁気抵抗は電気抵抗と同じ形なので、磁束が流れると熱が出る。}
\bkok{磁気抵抗はエネルギーをためるだけで、熱を出さない。鉄心の熱は、ヒステリシス損と渦電流損から出る。}
相似は「式の形」が同じということで、何でも同じではありません（⑤⑧）。
\end{tsumazuki}

\begin{tsumazuki}{すき間は小さいから、無視してよい？}
\bkng{\qty{1}{mm}のすき間は、\qty{20}{cm}の鉄心に比べて小さいので、無視してよい。}
\bkok{空気は鉄より何千倍も磁場を通しにくいので、\qty{1}{mm}でも磁気抵抗の9割以上をしめることがある。}
長さではなく、磁気抵抗 $l/(\mu S)$ で比べます（例題4）。
\end{tsumazuki}

\begin{tsumazuki}{アルミや銅は、磁場とまったく関係ない？}
\bkng{アルミや銅は磁石につかないので、磁場の中でも何も起きない。}
\bkok{アルミは常磁性、銅は反磁性で、とても弱い力を受ける。また、動かすと渦電流が流れて、強い力を受ける（電-8）。}
磁石につかないことと、磁場と関係ないことは違います（④）。
\end{tsumazuki}

%% ======================== 工学ではこう使う ========================
\Needspace{10\baselineskip}\bksection{phgear}{工学ではこう使う}
\begin{phkougaku}[鉄心と磁石を選ぶ]
\textbf{変圧器・モーターの鉄心}：鉄心には、保磁力が小さく（ヒステリシス損が小さい）、飽和が大きい電磁鋼板を使います。鉄心を薄い板に分けて重ねる（積層）のは、渦電流損を減らすためです（電-8 ⑪）。設計では、使う磁束密度を飽和（約\qty{2}{T}）の手前の1.5〜\qty{1.7}{T}くらいにし、磁気回路で要る巻数を決めます。

\textbf{磁気シールド}：病院の MRI の部屋や、精密な測定器を、外の磁場から守るために、透磁率の大きい材料で囲みます。磁力線は磁気抵抗の小さい道を通りたがるので、囲みの壁の中を回って、中の空間にはあまり入りません（磁気回路の並列：例題12 と同じ考え）。

\textbf{永久磁石}：モーターやスピーカーの磁石は、すき間の長さと磁石の長さの比で働く点が決まります（⑨）。磁石を長く（厚く）すると、すき間が同じでも磁場が強くなります。

\begin{phketa}[磁性材料の数の目安]
\phketaline{電磁鋼板の飽和}{約\qty{2}{T}}
\phketaline{電磁鋼板の鉄損（35A300）}{\qty{1.5}{T}・\qty{50}{Hz}で\qty{3.00}{W/kg}以下}
\phketaline{残留磁束密度}{ネオジム 約1.2〜\qty{1.4}{T}、フェライト 約\qty{0.4}{T}}
\phketaline{鉄のキュリー温度}{約770 ℃}
\phboxhead{答えが桁で合っているか}
鉄心の中の磁場が\qty{2}{T}をこえる答えは、飽和を無視しています。すき間のある電磁石なら、$B\approx\mu_0NI/l_{\mathrm g}$ で、「\qty{1}{mm}のすき間に\qty{1}{T}をつくるには約\qty{800}{A}ターン」と覚えておくと、すぐ見積もれます。
\end{phketa}
\end{phkougaku}
\phstation{電気機器,磁性材料,電力系統の機器,計測と磁気シールド}{透磁率、起磁力、磁気抵抗、磁気回路、鉄損、ヒステリシス損、保磁力、残留磁束密度、動作点、電磁鋼板}

%% ======================== 章末問題 ========================
\bkendproblems
\bkproblemgroup{基本}{3}
\bkproblem 1 m あたり2000回巻いたソレノイドに\qty{0.050}{A}を流し、比透磁率800の鉄心を入れます。中の磁場を求めましょう。
\bkpans{$H=\qty{100}{A/m}$。$B=800\times4\pi\times10^{-7}\times100$ で、約\qty{0.10}{T}}
\bkproblem 鉄心の中の磁場が\qty{1.6}{T}のとき、磁化はいくらですか（$H$ は小さいので無視する）。
\bkpans{$1.6\div(4\pi\times10^{-7})$ で、約 $1.3\times10^6$\,A/m}
\bkproblem 長さ\qty{0.40}{m}、断面積\qty{2.0}{cm^2}、比透磁率1500の鉄心に、300回巻いて\qty{0.20}{A}を流します。磁気抵抗、磁束、磁場を求めましょう。
\bkpans{$\mathcal R_{\mathrm m}\approx1.1\times10^6$\,/H、$\Phi=60\div(1.06\times10^6)$ で約 $5.7\times10^{-5}$\,Wb、$B\approx\qty{0.28}{T}$}
\bkproblem 問題3の鉄心に\qty{0.50}{mm}のすき間を作ると、磁場はいくらになりますか。
\bkpans{$\mathcal R_{\text{すき間}}\approx2.0\times10^6$\,/H。$\Phi=60\div(3.05\times10^6)$ で約 $2.0\times10^{-5}$\,Wb、$B\approx\qty{0.098}{T}$}
\bkproblem 地球の磁場の下向きの部分 $3.5\times10^{-5}$\,T（東京あたり。電-6 ⑫）を、磁場の強さ $H$ で表しましょう。
\bkpans{$3.5\times10^{-5}\div(4\pi\times10^{-7})$ で、約\qty{28}{A/m}}
\bkproblem 電気自動車のモーターの鉄心が、1周 \qty{200}{J/m^3} のループで\qty{400}{Hz}で磁化されています。1 kg あたりのヒステリシス損を求めましょう（密度\qty{7650}{kg/m^3}）。例題7（\qty{50}{Hz}）と比べて、何が言えますか。
\bkpans{$200\times400\div7650$ で、約\qty{10}{W/kg}。周波数が高いと、1周の損失が小さい材料でも、損失が大きくなる}
\bkproblemgroup{チャレンジ}{5}
\bkproblem 空気の中で、境目の法線と30°をなして、比透磁率2000の鉄へ入る磁力線は、鉄の中で法線と何度をなしますか。鉄の中の磁力線の向きについて、何が言えますか（例題8）。
\bkpans{$\tan\theta_{\text{鉄}}=2000\times\tan30^\circ\approx1150$ で、約89.95°。鉄の中では、ほとんど境目にそって進む}
\bkproblem ネオジム磁石（$B_{\mathrm r}=\qty{1.2}{T}$、$\mu_{\mathrm{rec}}=1.05$、長さ\qty{4.0}{mm}）が、\qty{2.0}{mm}のすき間に磁場をつくります。すき間の磁場と、すき間の面 1 cm$^2$ あたりの引き合う力（$B^2S/(2\mu_0)$）を求めましょう（例題11）。
\bkpans{$B=1.2\div(1+1.05\times0.50)$ で、約\qty{0.79}{T}。力は $0.787^2\times1.0\times10^{-4}\div(2\times4\pi\times10^{-7})$ で、約\qty{25}{N}}
\bkproblem 例題12の外鉄形の鉄心のまん中の脚に、長さ\qty{0.20}{mm}のすき間を作りました。まん中の磁束と磁場を求めましょう（例題4・12）。
\bkpans{$\mathcal R_{\text{すき間}}=l_{\mathrm g}/(\mu_0S)\approx4.0\times10^5$\,/H。全体 約 $8.0\times10^5$\,/H で、$\Phi\approx1.3\times10^{-4}$\,Wb、$B\approx\qty{0.31}{T}$}
\bkproblem コイルの巻数 $N$ と、1巻きを通る磁束 $\Phi$、電流 $I$ から、$L=N\Phi/I$ をコイルの自己インダクタンスといいます（電-8 ⑥）。例題3の鉄心と200回のコイルで、$L$ を磁気抵抗 $\mathcal R_{\mathrm m}$ で表して求めましょう。例題4の\qty{1.0}{mm}のすき間を作ると、$L$ は何分の1になりますか（例題3・4）。
\bkpans{$L=N^2/\mathcal R_{\mathrm m}$ で約\qty{0.20}{H}。すき間で約\qty{0.018}{H}、約11分の1}
\bkproblem 例題3のコイルの電流\qty{0.50}{A}を、\qty{0.010}{s}で一様に0まで減らします。コイルに生まれる電圧の大きさを、磁束の変化（ファラデーの法則：電-8 ③）から求めましょう。問題10の $L$ を使っても同じ答えになることを確かめましょう。同じ電流を\qty{0.0010}{s}で減らすと、電圧は何 V になりますか（電-8 ⑥・電-9 ④で、この電圧がスイッチに火花をとばすことを学びます）。
\bkpans{どちらも約\qty{10}{V}。\qty{0.0010}{s}なら約\qty{100}{V}}
\bkproblem 例題3の鉄心（すき間なし）の中は、$H=\qty{500}{A/m}$、$B\approx\qty{1.26}{T}$ です。鉄心の中の磁場のエネルギー密度を $\tfrac12BH$ として、鉄心全体（体積は長さ×断面積）にたまるエネルギーを求めましょう。問題10の $L$ を使って $\tfrac12LI^2$（電-8 ⑧）を求め、比べましょう。
\bkpans{$\tfrac12\times1.26\times500\approx\qty{314}{J/m^3}$、体積 $8.0\times10^{-5}$\,m$^3$ で約 $2.5\times10^{-2}$\,J。$\tfrac12\times0.20\times0.50^2$ も約 $2.5\times10^{-2}$\,J で、同じになる}

\begin{furikaeri}
\begin{checklist}
\item 磁石に近づけた鉄が磁石になるわけを、小さな磁石の向きがそろう図で説明できる。\bkfurref{ph:電-7:sec2}
\item 磁化 $M$・磁場の強さ $H$・比透磁率を使って、鉄心の中の磁場を求められる。\bkfurref{ph:電-7:sec3}
\item 磁気回路を電気回路と同じ形で解き、すき間のある電磁石の磁束を求められる。\bkfurref{ph:電-7:sec5}
\item 鉄と空気の境目で磁力線がどう曲がるかを、境界条件から求められる。\bkfurref{ph:電-7:sec7}
\item $BH$ 曲線の面積からヒステリシス損を、減磁曲線から永久磁石の動作点を求められる。\bkfurref{ph:電-7:sec8}
\end{checklist}
\end{furikaeri}
\bkfukidashi{mina}{くぎの中の小さな磁石がそろうから、くぎも磁石になったんだね。鉄心の磁気回路が、電気回路と同じ形で解けるなんて、おもしろい！}
